\documentclass{article}

\PassOptionsToPackage{numbers, compress}{natbib}
% if you need to pass options to natbib, use, e.g.:
%     \PassOptionsToPackage{numbers, compress}{natbib}
% before loading neurips_2025


% ready for submission
\usepackage{neurips_2025}


% to compile a preprint version, e.g., for submission to arXiv, add add the
% [preprint] option:
%     \usepackage[preprint]{neurips_2025}


% to compile a camera-ready version, add the [final] option, e.g.:
%     \usepackage[final]{neurips_2025}


% to avoid loading the natbib package, add option nonatbib:
%    \usepackage[nonatbib]{neurips_2025}


\usepackage[utf8]{inputenc} % allow utf-8 input
\usepackage[T1]{fontenc}    % use 8-bit T1 fonts
\usepackage{hyperref}       % hyperlinks
\usepackage{url}            % simple URL typesetting
\usepackage{booktabs}       % professional-quality tables
\usepackage{amsfonts}       % blackboard math symbols
\usepackage{mathrsfs}
\usepackage{amsthm,amssymb,amsmath}
\newtheorem{proposition}{Proposition}
\newtheorem{definition}{Definition}  % 定义“定义”环境
\newtheorem{theorem}{Theorem}
\newtheorem{lemma}{Lemma}



\usepackage{nicefrac}       % compact symbols for 1/2, etc.
\usepackage{microtype}      % microtypography
\usepackage{xcolor}         % colors
\usepackage{graphicx}
\usepackage{float}
\usepackage{algorithm}
\usepackage{algorithmic}

\usepackage{xcolor}
\usepackage{tabularx}


\usepackage{booktabs,tabularx,xcolor,array}
% Optional: 
\usepackage{caption}
% Color macros for per-method annotation
\newcommand{\colorred}[1]{\textcolor{red!70!black}{#1}}      % soft penalty
\newcommand{\colorblue}[1]{\textcolor{blue!70!black}{#1}}    % hard projection
\newcommand{\colorgreen}[1]{\textcolor{green!45!black}{#1}}  % architectural isolation
\newcommand{\colorgray}[1]{\textcolor{black!55}{#1}}         % none / no explicit constraint

\title{Bridging Regularization and Flatness in Continual Learning via Drift and Diffusion Control}

%Regularization Meets Flatness in Continual Learning Through Curvature–Noise Alignment

% The \author macro works with any number of authors. There are two commands
% used to separate the names and addresses of multiple authors: \And and \AND.
%
% Using \And between authors leaves it to LaTeX to determine where to break the
% lines. Using \AND forces a line break at that point. So, if LaTeX puts 3 of 4
% authors names on the first line, and the last on the second line, try using
% \AND instead of \And before the third author name.


\author{%
  David S.~Hippocampus\thanks{Use footnote for providing further information
    about author (webpage, alternative address)\emph{not} for acknowledging
    funding agencies.} \\
  Department of Computer Science\\
  Cranberry-Lemon University\\
  Pittsburgh, PA 15213 \\
  \texttt{hippo@cs.cranberry-lemon.edu} \\
  % examples of more authors
  % \And
  % Coauthor \\
  % Affiliation \\
  % Address \\
  % \texttt{email} \\
  % \AND
  % Coauthor \\
  % Affiliation \\
  % Address \\
  % \texttt{email} \\
  % \And
  % Coauthor \\
  % Affiliation \\
  % Address \\
  % \texttt{email} \\
  % \And
  % Coauthor \\
  % Affiliation \\
  % Address \\
  % \texttt{email} \\
}


\begin{document}


\maketitle


\begin{abstract}
Most continual learning approaches mitigate catastrophic forgetting by regulating only the deterministic component of optimization---modifying gradient directions to reduce drift interference with previous tasks. However, stochastic training dynamics introduce a second, largely overlooked mechanism: the interaction between optimization noise and the curvature of previously learned tasks. We provide a theoretical framework, grounded in It\^{o} calculus, that decomposes the expected forgetting into two terms: (i)~\emph{drift interference}, governed by the alignment of old-task gradients with the current update direction, and (ii)~\emph{diffusion interference}, governed by the cross-task noise--curvature term $\mathrm{tr}(H_1\Sigma)$. This decomposition reveals that both regularization-based and gradient projection methods address only drift, leaving diffusion interference uncontrolled. We formulate continual learning as a constrained optimization over both drift and diffusion, and derive a structured noise covariance $\Sigma^* \propto (F_1 - \beta F_2 + \varepsilon I)^{-1}$ that suppresses perturbations along high-curvature directions of old tasks while preserving exploration for the current task. Experiments on ImageNet-R confirm the theoretical predictions: isotropic Gaussian noise injection significantly amplifies forgetting ($-16.7$ AAA points vs.\ OGD), Fisher-guided diffusion control alone improves OGD by $+8.2$ AAA points, and the combined drift--diffusion controller achieves AAA$=70.3$ ($+29.6$ over the OGD baseline), validating the joint control objective.
\end{abstract}


% \section{Notions}

% This section collects the main symbols used throughout the paper and follows the order in which they first appear.

% \subsection{Preliminaries (information theory, martingales, and subspace geometry)}
% \begin{itemize}
% \item $(\Omega,\mathcal F)$: measurable space.
% \item $\mathbb P(\cdot)$, $\mathbb E[\cdot]$: probability and expectation.
% \item $Q\ll P$: absolute continuity of measures.
% \item $\frac{dQ}{dP}$: Radon--Nikodym derivative.
% \item $p,q$: densities of $P,Q$ with respect to a common dominating measure.
% \item $\otimes$: product measure.
% \item $P_{UV}$, $P_U$, $P_V$: joint and marginal distributions.
% \item $P_{V\mid U}$: conditional distribution.
% \item $\mathrm{KL}(Q\|P)$: Kullback--Leibler divergence between probability measures $Q$ and $P$.
% \item $I(U;V)$: mutual information between random variables $U$ and $V$.
% \item $\mathbb F=\{\mathcal F_t\}_{t\ge 0}$: a filtration (increasing family of $\sigma$-algebras).
% \item $X_t$: a nonnegative adapted process.
% \item $M_t$: a (super)martingale adapted to $\mathbb F$.
% \item $A=U\Sigma V^\top$: singular value decomposition (SVD) of a matrix $A$.
% \item $\sigma_i(A)$: the $i$-th singular value of $A$.
% \item $\mathrm{Gr}(k,D)$: Grassmann manifold of $k$-dimensional subspaces of $\mathbb R^D$.
% \item $\mathcal U,\mathcal V\subseteq\mathbb R^D$: $k$-dimensional subspaces.
% \item $U,V\in\mathbb R^{D\times k}$: column-orthonormal basis matrices spanning $\mathcal U$ and $\mathcal V$.
% \item $\{\theta_i\}_{i=1}^k$: principal angles between $\mathcal U$ and $\mathcal V$.
% \item $d_{\mathrm{geo}}(\mathcal U,\mathcal V)$, $d_{\mathrm{chord}}(\mathcal U,\mathcal V)$: geodesic and chordal distances on $\mathrm{Gr}(k,D)$.
% \end{itemize}

% \subsection{Supervised learning}
% \begin{itemize}
% \item $\mathcal X$: input space.
% \item $\mathcal Y$: label/output space.
% \item $\mathcal Z:=\mathcal X\times\mathcal Y$: sample space, with $z=(x,y)$.
% \item $\mathcal D$: data distribution over $\mathcal Z$.
% \item $S=\{z_j\}_{j=1}^m$: training set drawn i.i.d. from $\mathcal D$.
% \item $m:=|S|$: sample size.
% \item $W\in\mathbb R^d$: model parameter vector (dimension $d$).
% \item $f_W:\mathcal X\to\mathcal Y$: predictor induced by parameters $W$.
% \item $\ell(W,z)$: loss.
% \item $L_{\mathcal D}(W):=\mathbb E_{z\sim\mathcal D}\,\ell(W,z)$: population risk.
% \item $\widehat L_{S}(W):=\frac{1}{m}\sum_{j=1}^{m}\ell(W,z_j)$: empirical risk.
% \item $\widehat W$: a particular solution/iterate learned on $S$.
% \end{itemize}

% \subsection{Spectral representation}
% \begin{itemize}
% \item $\phi_W(x)$: representation map.
% \item $h=\phi_W(x)$: representation vector.
% \item $\Sigma_h=\mathbb E[h h^\top]$: second moment / covariance (centered if needed).
% \item $\lambda_i(\Sigma_h)$: $i$-th eigenvalue of $\Sigma_h$.
% \end{itemize}

% \subsection{Continual learning and PEFT/LoRA}
% \begin{itemize}
% \item $T$: number of tasks.
% \item $t\in\{1,\dots,T\}$: task index.
% \item $\mathcal D_t$: data distribution of task $t$.
% \item $S_t=\{z_{t,j}\}_{j=1}^{m_t}$: training set for task $t$ drawn i.i.d. from $\mathcal D_t$.
% \item $m_t:=|S_t|$: number of training samples in task $t$.
% \item $W_0$: pretrained initialization.
% \item $W_t$: parameters after learning tasks $1{:}t$.
% \item $W_t^{\mathrm{froz}}$: frozen component when learning task $t$.
% \item $\Delta W_t$: trainable increment at task $t$.
% \item $W_t=W_t^{\mathrm{froz}}+\Delta W_t$: frozen-plus-increment decomposition.
% \item $\widehat{\Delta W}_t$: learned increment at task $t$.
% \item $\widehat W_t:=W_t^{\mathrm{froz}}+\widehat{\Delta W}_t$.
% \item $\mathcal W_{\Delta,t}\subseteq\mathbb R^d$: feasible update subspace at task $t$.
% \item $r$: LoRA rank.
% \item $A\in\mathbb R^{r\times d_{\mathrm{in}}},\;B\in\mathbb R^{d_{\mathrm{out}}\times r}$: LoRA factors such that $\Delta W=BA$.
% \item $\alpha>0$: (optional) LoRA scaling parameter, often used as $\Delta W=\frac{\alpha}{r}BA$.
% \end{itemize}

% \subsection{PAC-Bayesian theorem}
% \begin{itemize}
% \item $\mathcal H$: hypothesis space.
% \item $\mathcal M(\mathcal H)$: set of probability measures over $\mathcal H$.
% \item $P\in\mathcal M(\mathcal H)$: prior distribution.
% \item $Q(\cdot\mid P,S)\in\mathcal M(\mathcal H)$: posterior distribution.
% \item $\mathrm{KL}(Q\|P)$: KL divergence term appearing in PAC-Bayes bounds.
% \item $\delta\in(0,1)$: confidence level parameter.

% \item $\mathcal P_t\in\mathcal M\big(\mathcal M(\mathcal H)\big)$: a \emph{hyperdistribution} (a probability measure over priors $P\in\mathcal M(\mathcal H)$).
% \item $\mathcal P_0$: hyperprior.
% \item $\mathcal P_t$: hyperposterior after observing data up to time $t$.
% \item $P_t\sim\mathcal P_{t-1}$: a sampled task-dependent prior drawn from the hyperdistribution.
% \item $Q_t(\cdot\mid P_t,S_t)$: task-$t$ posterior obtained from prior $P_t$ and data $S_t$.

% \item $L_{\mathcal D}(Q)=\mathbb E_{W\sim Q}\,L_{\mathcal D}(W)$: posterior-averaged population risk.
% \item $\widehat L_{S}(Q)=\mathbb E_{W\sim Q}\,\widehat L_{S}(W)$: posterior-averaged empirical risk.

% \item $\mathcal L_t:=\mathbb E_{P_t\sim\mathcal P_{t-1}}\,\mathbb E_{W_t\sim Q_t(\cdot\mid P_t,S_t)}\big[L_{\mathcal D_t}(W_t)\big]$: hierarchical (hyper) test risk.
% \item $\widehat{\mathcal L}_t:=\mathbb E_{P_t\sim\mathcal P_{t-1}}\,\mathbb E_{W_t\sim Q_t(\cdot\mid P_t,S_t)}\big[\widehat L_{S_t}(W_t)\big]$: hierarchical (hyper) empirical risk.
% \end{itemize}

% \subsection{Sharpness and curvature}
% \begin{itemize}
% \item $\|\cdot\|$: norm on parameter space (default: Euclidean norm).
% \item $\varepsilon$ (or $\epsilon$): parameter perturbation.
% \item $\rho>0$: neighborhood radius.
% \item $\Sigma\succeq 0$: covariance of (e.g., Gaussian) perturbation.
% \item $\nabla_W\widehat L_S(W)$: empirical-risk gradient.
% \item $H_S(W):=\nabla_W^2\widehat L_S(W)$: empirical-risk Hessian.
% \item $\Pi_{\Delta,t}$: projector onto the LoRA update subspace $\mathcal W_{\Delta,t}$.
% \item $C_t := H_{S_t}(\widehat W_t)=\nabla_W^2\widehat L_{S_t}(\widehat W_t)$: task-$t$ local curvature.
% \item $C_{\Delta,t} := \Pi_{\Delta,t} C_t \Pi_{\Delta,t}$: curvature restricted to $\mathcal W_{\Delta,t}$.
% \item $\operatorname{tr}(\cdot)$: matrix trace.
% \item $\mathrm{AS}^{(0)}_S(W,\rho)$: adversarial zeroth-order sharpness.
% \item $\mathrm{RS}^{(0)}_S(W,\Sigma)$: random zeroth-order sharpness.
% \item $\mathrm{AS}^{(1)}_S(W,\rho)$: adversarial first-order sharpness.
% \item $\kappa_{\mathrm{Tr}}^{\phi}(\mathbf w)$: relative flatness under feature map $\phi$.
% \item $\mathbf w\in\mathbb R^{d\times m}$: linear predictor weights.
% \item $H_{s,s'}(\mathbf w,\phi(S))$: block Hessian of the empirical risk.
% \end{itemize}

% \section{Preliminaries}

% \subsection{Kullback--Leibler (KL) divergence}
% Let $(\Omega,\mathcal F)$ be a measurable space and let $P,Q$ be probability measures on it.
% If $Q\ll P$ (i.e., $Q$ is absolutely continuous with respect to $P$), the \emph{Kullback--Leibler divergence} from $Q$ to $P$ is defined by
% \begin{equation}
% \mathrm{KL}(Q\|P)
% :=
% \int \log\Big(\frac{dQ}{dP}\Big)\,dQ.
% \end{equation}
% If $P,Q$ admit densities $p,q$ with respect to a common dominating measure (e.g., Lebesgue), then
% \begin{equation*}
% \mathrm{KL}(Q\|P)
% =
% \int q(x)\,\log\frac{q(x)}{p(x)}\,dx.
% \end{equation*}
% We use the convention $\mathrm{KL}(Q\|P)=+\infty$ when $Q\not\ll P$.

% \subsection{Mutual information}
% Let $U,V$ be random variables on a common probability space with joint distribution $P_{UV}$ and marginals $P_U,P_V$.
% In particular, $I(U;V)\ge 0$, and $I(U;V)=0$ holds if and only if $P_{UV}=P_U\otimes P_V$, i.e., $U$ and $V$ are independent.
% The \emph{mutual information} between $U$ and $V$ is defined as the KL divergence between the joint law and the product of marginals:
% \begin{equation}
% I(U;V)
% :=
% \mathrm{KL}\big(P_{UV}\,\|\,P_U\otimes P_V\big).
% \end{equation}
% If $(U,V)$ admits a joint density $p(u,v)$ with marginals $p(u),p(v)$, then
% \begin{equation*}
% I(U;V)
% =
% \int p(u,v)\,\log\frac{p(u,v)}{p(u)p(v)}\,dudv
% =
% \mathbb E\Big[\log\frac{p(U,V)}{p(U)p(V)}\Big].
% \end{equation*}
% Equivalently, mutual information can be expressed as an expected conditional KL divergence:
% \begin{equation*}
% I(U;V)
% =
% \int p(u)\,\mathrm{KL}\big(P_{V\mid U=u}\,\|\,P_V\big)\,du
% =
% \mathbb E_{U}\Big[\mathrm{KL}\big(P_{V\mid U}\,\|\,P_V\big)\Big],
% \end{equation*}
% which makes explicit that $I(U;V)$ measures, on average, how far the conditional distribution $P_{V\mid U}$ deviates from the marginal $P_V$.

% % \paragraph{Relationship to KL divergence and dependence.}
% % By definition, mutual information is a KL divergence:
% % \begin{equation}
% % I(U;V)=\mathrm{KL}\big(P_{UV}\,\|\,P_U\otimes P_V\big).
% % \end{equation}
% % Consequently, it inherits the standard properties of KL divergence.
% % In particular, $I(U;V)\ge 0$, and $I(U;V)=0$ holds if and only if $P_{UV}=P_U\otimes P_V$, i.e., $U$ and $V$ are independent.
% % Moreover, writing the joint density as $p(u,v)=p(u)\,p(v\mid u)$ yields
% % \begin{equation}
% % I(U;V)
% % =
% % \int p(u,v)\,\log\frac{p(v\mid u)}{p(v)}\,dudv
% % =
% % \mathbb E\Big[\log\frac{p(V\mid U)}{p(V)}\Big].
% % \end{equation}
% % Equivalently, mutual information can be expressed as an expected conditional KL divergence:
% % \begin{equation}
% % I(U;V)
% % =
% % \int p(u)\,\mathrm{KL}\big(P_{V\mid U=u}\,\|\,P_V\big)\,du
% % =
% % \mathbb E_{U}\Big[\mathrm{KL}\big(P_{V\mid U}\,\|\,P_V\big)\Big],
% % \end{equation}
% % which makes explicit that $I(U;V)$ measures, on average, how far the conditional distribution $P_{V\mid U}$ deviates from the marginal $P_V$.



% \subsection{Filtration and martingales}

% \paragraph{Filtration (information).}
% Let $(\Omega,\mathcal F,\mathbb P)$ be a probability space and let $I$ be a totally ordered index set (e.g., $I=\{0,1,2,\dots\}$). A \emph{filtration} is a family of sub-$\sigma$-algebras $\mathbb F:=\{\mathcal F_t\}_{t\in I}$ such that
% \begin{equation*}
% \mathcal F_s\subseteq \mathcal F_t\subseteq \mathcal F,\qquad \forall\, s\le t.
% \end{equation*}
% Intuitively, $\mathcal F_t$ represents the information available up to time $t$. In our task-sequence setting, a natural choice is
% \begin{equation*}
% \mathcal F_t:=\sigma(S_1,\dots,S_t),
% \end{equation*}
% so $\mathcal F_{t-1}=\sigma(S_1,\dots,S_{t-1})$ is the information available before observing task $t$.

% \paragraph{Adaptedness / measurability.}
% A process $\{X_t\}_{t\in I}$ is \emph{adapted} to $\mathbb F$ if each $X_t$ is $\mathcal F_t$-measurable. This formalizes the idea that $X_t$ can be determined using only the information available by time $t$.

% \paragraph{Martingale and supermartingale.}
% Let $\{M_t\}_{t\in I}$ be an integrable process (i.e., $\mathbb E[|M_t|]<\infty$ for all $t$) adapted to $\mathbb F$.
% \begin{itemize}
% \item $\{M_t\}$ is a \emph{martingale} (w.r.t. $\mathbb F$) if for all $t\ge 1$,
% \begin{equation*}
% \mathbb E[M_t\mid \mathcal F_{t-1}] = M_{t-1} \quad \text{a.s.}
% \end{equation*}
% \item $\{M_t\}$ is a \emph{supermartingale} (w.r.t. $\mathbb F$) if for all $t\ge 1$,
% \begin{equation*}
% \mathbb E[M_t\mid \mathcal F_{t-1}] \le M_{t-1} \quad \text{a.s.}
% \end{equation*}
% (Analogously, $\ge$ defines a submartingale.)
% \end{itemize}

% \paragraph{Product supermartingale induced by conditional mean control.}
% Assume $X_t\ge 0$, $X_t$ is $\mathcal F_t$-measurable, and
% \begin{equation*}
% \mathbb E[X_t\mid \mathcal F_{t-1}]\le 1.
% \end{equation*}
% Define
% \begin{equation*}
% M_0:=1,\qquad M_t:=\prod_{s=1}^t X_s.
% \end{equation*}
% Then $\{M_t\}_{t\ge 0}$ is a nonnegative supermartingale since
% \begin{equation*}
% \mathbb E[M_t\mid \mathcal F_{t-1}]
% = M_{t-1}\,\mathbb E[X_t\mid \mathcal F_{t-1}]\le M_{t-1}.
% \end{equation*}
% In particular, taking expectations yields $\mathbb E[M_T]\le \mathbb E[M_0]=1$. For any $a>0$, Markov's inequality gives $\mathbb P(M_T\ge a)\le \mathbb E[M_T]/a\le 1/a$.

% \subsection{Subspace distance}

% \paragraph{SVD and singular values.}
% For any matrix $A\in\mathbb R^{m\times n}$, the singular value decomposition (SVD) writes
% \begin{equation}
% A = U\,\Sigma\,V^\top,
% \end{equation}
% where $U\in\mathbb R^{m\times m}$ and $V\in\mathbb R^{n\times n}$ are orthogonal, and $\Sigma\in\mathbb R^{m\times n}$ is diagonal (rectangular) with nonnegative diagonal entries $\sigma_1(A)\ge\sigma_2(A)\ge\cdots\ge 0$ called the \emph{singular values}.
% Equivalently, $\sigma_i(A)=\sqrt{\lambda_i(A^\top A)}=\sqrt{\lambda_i(AA^\top)}$.

% \paragraph{Principal angles between subspaces (SVD characterization).}
% Let $\mathcal U,\mathcal V\subseteq\mathbb R^D$ be two $k$-dimensional subspaces and let $U,V\in\mathbb R^{D\times k}$ be column-orthonormal basis matrices with $\mathrm{span}(U)=\mathcal U$ and $\mathrm{span}(V)=\mathcal V$.
% The \emph{principal angles} $0\le\theta_1\le\cdots\le\theta_k\le \frac{\pi}{2}$ are defined via the SVD
% \begin{equation}
% U^\top V = P\,\mathrm{diag}(\cos\theta_1,\dots,\cos\theta_k)\,Q^\top,
% \end{equation}
% so that $\cos\theta_i=\sigma_i(U^\top V)$; see, e.g.,~\citep{hammGrassmannDiscriminantAnalysis2008}.

% \paragraph{Common subspace distance metrics.}
% The principal angles induce several standard distances on the Grassmann manifold $\mathrm{Gr}(k,D)$.
% Two widely used ones are:
% \begin{itemize}
% \item \emph{Chordal / projection distance:}
% \begin{equation}
%  d_{\mathrm{chord}}(\mathcal U,\mathcal V)
% :=\|\sin\theta\|_2
% =\Big(\sum_{i=1}^k \sin^2\theta_i\Big)^{1/2}
% =\frac{1}{\sqrt{2}}\,\|UU^\top - VV^\top\|_F.
% \end{equation}
% \item \emph{Geodesic distance:}
% \begin{equation*}
%  d_{\mathrm{geo}}(\mathcal U,\mathcal V)
% :=\|\theta\|_2
% =\Big(\sum_{i=1}^k \theta_i^2\Big)^{1/2}.
% \end{equation*}

% \end{itemize}
% Both distances depend only on the subspaces (not on the specific choice of orthonormal bases).

% \section{Basic Setting}

% \begin{figure}[h]
%   \centering
%   \includegraphics[width=0.8\textwidth]{figs/FeatureMatrix.drawio.pdf}
%   \caption{Decompes of the model.}
%   \label{fig:example_image}
% \end{figure}


% \subsection{Supervised learning}

% % \paragraph{Spaces and model factorization.}
% Let $\mathcal X$ be the input space, $\mathcal Y$ the label space, and $\mathcal Z=\mathcal X\times\mathcal Y$ the data space.
% A model is parameterized by $W=(W_\phi,W_c)$, where $W_\phi$ denotes the parameters of the feature extractor and $W_c$ denotes the parameters of the readout head.
% The induced predictor is written as $f_W:\mathcal X\to\mathcal Y$.

% The feature extractor induces a representation map
% \begin{equation}
% \label{eq: def feature}
% \phi_{W_\phi}:\mathcal X\to\mathbb R^{p},
% \qquad
% h=\phi_{W_\phi}(x)\in\mathbb R^{p}.
% \end{equation}
% The readout head acts on the representation space, so that the overall predictor factorizes as
% \begin{equation}
% f_W(x) = W_c\bigl(\phi_{W_\phi}(x)\bigr).
% \end{equation}


% \textbf{Spectral view}

% \textbf{Inner-product readout and label embeddings.}
% For multi-class classification, we use the standard linear-logit parameterization as a concrete instantiation of the head.
% Specifically, we define a scoring function
% \begin{equation*}
% s_W(x,y)
% :=
% \bigl\langle \phi_{W_\phi}(x),\ \mu_{W_c}(y)\bigr\rangle,
% \qquad
% \mu_{W_c}:\mathcal Y\to\mathbb R^{p},
% \end{equation*}
% where $\mu_{W_c}(y)$ denotes the label embedding induced by the head parameters (for a linear classifier, $\mu_{W_c}(y)$ is the weight vector of class $y$).
% The predictive distribution is then given by
% \begin{equation*}
% p_W(y\mid x)=\frac{\exp(s_W(x,y))}{\sum_{y'\in\mathcal Y}\exp(s_W(x,y'))}.
% \end{equation*}
% This form makes explicit that, at the level of logits, the model compares an input representation $\phi_{W_\phi}(x)$ with a label-side embedding $\mu_{W_c}(y)$ via an inner product.


% \textbf{conditional operators and sufficient representations.}
% Beyond model parameterization, it is useful to view learning as approximating the conditional operator $P(\cdot\mid\cdot)$.
% For any conditional distribution $P(y\mid x)$, one may express it as a (possibly infinite-dimensional) singular-function expansion
% \begin{equation*}
% P(y\mid x)=\bigl\langle \varphi(x),\ \mu(y)\bigr\rangle,
% \end{equation*}
% where $\varphi:\mathcal X\to\mathbb R^{d}$ and $\mu:\mathcal Y\to\mathbb R^{d}$ are the input-side and output-side singular embeddings, respectively, and $d$ can be infinite.
% In this viewpoint, a representation is \emph{sufficient} for downstream prediction when it preserves the dominant spectral components of the conditional operator, so that task-relevant information remains recoverable by simple readouts.
% Our parameterized model can be interpreted as a finite-dimensional approximation of such a factorization, with $\phi_{W_\phi}(x)$ serving as a learned input embedding and $\mu_{W_c}(y)$ serving as a learned label embedding.

% \paragraph{Loss.}
% We consider a bounded loss $\ell:\mathcal H\times\mathcal Z\to[a,b]$ and write $\ell(W,z)$ as shorthand for $\ell(f_W,z)$, where $\mathcal H$ denotes the hypothesis class induced by the chosen parameterization.























% % Let $\mathcal X$ be the input space, $\mathcal Y$ the label space, and $\mathcal Z = \mathcal X \times \mathcal Y$ the data space.

% % A model is parameterized by
% % $
% % W=(W,W_c)\in\mathbb R^d,
% % $
% % where $W$ denotes the parameters of the feature extractor (representation encoder) and $W_c$ denotes the parameters of the classifier (readout head). The induced predictor is written as
% % $
% % f_W:\mathcal X\to\mathcal Y.
% % $

% % The feature extractor with parameters $W$ induces a representation map
% % $
% % \phi_{W}:\mathcal X\to\mathbb R^p,
% % \qquad
% % h=\phi_{W}(x)\in\mathbb R^p.
% % $
% % The classifier with parameters $W_c$ acts on the representation space via a readout map
% % $
% % {W_c}:\mathbb R^p\to\mathcal Y.
% % $
% % Consequently, the overall predictor factorizes as the composition
% % $
% % f_W(x)={W_c} \bigl(\phi_{W}(x)\bigr).
% % $


% % We consider a bounded loss $\ell : \mathcal H \times \mathcal Z \to [0,1]$ and write $\ell(W,z)$ for $\ell(f_W,z)$. (The bound is not necessary limited to $[0,1]$, it is can easily be extended to [a,b]).


% Given a distribution $\mathcal D$ over $\mathcal Z$ and a training set $S=\{z_j\}_{j=1}^m$ drawn i.i.d.\ from $\mathcal D$, the population and empirical risks are
% \begin{equation}
% L_{\mathcal D}(W)
% =
% \mathbb E_{z\sim\mathcal D}\,\ell(W,z),
% \qquad
% \widehat L_{S}(W)
% =
% \frac{1}{m}\sum_{j=1}^{m}\ell(W,z_{j}).
% \end{equation}

% \subsection{Spectral Representation}
% This subsection states the minimal notation for spectral analysis of representations.

% % \paragraph{Representation map.}
% % Given a model/encoder with parameters $W$, we denote its representation map by
% % $
% % \phi_{W}:\mathcal X\to\mathbb R^p,
% % \qquad
% % h=\phi_{W}(x)\in\mathbb R^p.
% % $

% \paragraph{Second moment.}
% Based on eq \ref{eq: def feature}.
% For $x\sim\mathcal D$, define the (uncentered) second moment (or centered covariance, if needed) as
% \begin{equation}
% \Sigma_h := \mathbb E_{x\sim\mathcal D}\big[h h^\top\big] \in \mathbb R^{p\times p}.
% \end{equation}
% We write its spectral decomposition as $\Sigma_h = U\Lambda U^\top$, where $\Lambda=\mathrm{diag}(\lambda_1,\dots,\lambda_p)$ and $\lambda_1\ge\cdots\ge\lambda_p\ge 0$.


% Given a sample $S=\{(x_j,y_j)\}_{j=1}^m$, let $h_j=\phi_W(x_j)$ and define the feature matrix
% \begin{equation}
% \label{eq:feature matrix}
% H := \begin{bmatrix} h_1^\top \\ \vdots \\ h_m^\top \end{bmatrix}\in\mathbb R^{m\times p}.
% \end{equation}

% \paragraph{the empirical second moment of feature matrix}
% Then the empirical second moment is 
% \begin{equation}
% \label{eq:the empirical second moment of feature matrix}
%   \widehat\Sigma_h := \frac{1}{m}H^\top H.
% \end{equation}


% % \paragraph{Singular values and principal angles ().}
% For any matrix $A$, denote its singular values by $\{\sigma_i(A)\}_i$. In particular,
% $\lambda_i(H^\top H)=\sigma_i(H)^2$.



% \subsection{LoRA parameterization}
% We denote the pretrained (frozen) parameters by $W_0$.
% LoRA introduces a trainable low-rank update $\Delta W$ and performs adaptation via
% \begin{equation}
% W = W_0 + \Delta W.
% \end{equation}
% % For a weight matrix $W_0\in\mathbb R^{d_{\mathrm{out}}\times d_{\mathrm{in}}}$, LoRA parameterizes the update as
% % \begin{equation}
% % \Delta W := B A,
% % \qquad
% % B\in\mathbb R^{d_{\mathrm{out}}\times r},\;A\in\mathbb R^{r\times d_{\mathrm{in}}},\; r\ll \min\{d_{\mathrm{out}},d_{\mathrm{in}}\},
% % \end{equation}
% % so that $\mathrm{rank}(\Delta W)\le r$.
% % A common scaled variant is $\Delta W := \frac{\alpha}{r}BA$ for some scalar $\alpha>0$.
% % In all cases, the pretrained weights $W_0$ are kept fixed and only $A,B$ (equivalently $\Delta W$ in the induced low-dimensional subspace) are optimized.
% Concretely, for a target matrix
% $
% W_{0}^{(\ell)}\in\mathbb R^{d_{\mathrm{out}}\times d_{\mathrm{in}}},
% $
% LoRA parameterizes the update as
% $
% W^{(\ell)} = W_{0}^{(\ell)} + \Delta W^{(\ell)},
% \qquad
% \Delta W^{(\ell)} = B^{(\ell)}A^{(\ell)},
% $
% with
% $
% B^{(\ell)}\in\mathbb R^{d_{\mathrm{out}}\times r},
% \quad
% A^{(\ell)}\in\mathbb R^{r\times d_{\mathrm{in}}},
% \quad
% r\ll \min{d_{\mathrm{out}},d_{\mathrm{in}}},
% $
% so that $\mathrm{rank}(\Delta W_{\phi}^{(\ell)})\le r$. A common scaled variant is
% $
% \Delta W^{(\ell)} := \frac{\alpha}{r}B^{(\ell)}A^{(\ell)}
% $
% for some scalar $\alpha>0$.
% All pretrained backbone weights $W_{0}$ are kept fixed, and only the LoRA factors ${A^{(\ell)},B^{(\ell)}}_\ell$ are optimized.


% \textbf{Training a new classifier head with LoRA}
% We attach and train a task-specific classifier head $W_c$ on top of the adapted representation. In particular, we treat $W_c$ as fully trainable parameters (initialized from scratch) and optimize them jointly with the LoRA factors in $W$, while keeping the pretrained backbone weights $W_{0}$ frozen

% \subsection{Continual learning (task sequence) }
% We consider a sequence of tasks indexed by $t\in\{1,\dots,T\}$. Each task $t$ is associated with a data distribution $\mathcal D_t$ over $\mathcal Z$ and a training set
% $S_t=\{z_{t,j}\}_{j=1}^{m_t}$ of size $m_t:=|S_t|$, where samples are drawn i.i.d.\ from $\mathcal D_t$.
% For any parameter vector $W$, we define the task-$t$ population/empirical risks by
% \begin{equation}
% L_{\mathcal D_t}(W)=\mathbb E_{z\sim\mathcal D_t}\,\ell(W,z),
% \qquad
% \widehat L_{S_t}(W)=\frac{1}{m_t}\sum_{z\in S_t}\ell(W,z).
% \end{equation}


% 保留参数
% \begin{algorithm}[H]
% \caption{Algorithm}\label{alg:alg1}
% \begin{algorithmic}
% \STATE {\textsc{Input} Dataset for tasks: $S_T$, task: $T$, initial model: $\theta_0$, reply  Memory Buffer: $M$
% }
% \State {\textsc{Cache between task} Memory Buffer: $M$,  Model after trained: $\theta_t$,
% }
% \State {\textsc{Output} final model $\theta_T$
% }
% \STATE  \textbf{First task} $t=1$
% \STATE \hspace{0.5cm} {init timely memory buffer} $M_0$
% \STATE \hspace{0.5cm} \textbf{Batch} $i=1,\dots,N$
% \STATE \hspace{1cm} \textbf{Sample batch} Sample batch $B_i$ from $S_1$
% \STATE \hspace{1cm} \textbf{Compute gradient} batch’s loss gradient using optimizer  $g_{S_1,B_i}=\nabla L_{S_1,B_i}(\theta)$
% \STATE \hspace{1.5cm} \textbf{SGD} $g_{S,B_i}=\nabla L_{S,B_i}(\theta)$

% \STATE \hspace{1.5cm} \textbf{RWP} 
% \STATE \hspace{2cm} \textbf{zero-order random perturbation} $\delta^0 \sim \mathcal{N}(0,\rho I)$
% \STATE \hspace{2cm}$g_{S,B_i}=\nabla L_{S,B_i}(\theta + \delta^0)$

% \STATE \hspace{1.5cm} \textbf{SAM}
% \STATE \hspace{2cm} \textbf{zero-order worst perturbation} $\delta^0 = \rho  \cdot \frac{\nabla L_{S,B_i}(\theta)}{\| \nabla L_{S,B_i}(\theta)\|_2}$
% \STATE \hspace{2cm} $g_{S,B_i}= \nabla L_{S,B_i}(\theta+\delta^0)$

% \STATE \hspace{1.5cm} \textbf{GAM}: 
% % $g_{S_1,B_i}=\nabla L_{S_1,B_i}(\theta ) + \alpha g_{norm}$ 
% \STATE \hspace{2cm} \textbf{first-order worst perturbation} $\delta^1=\rho \frac{\nabla\|\nabla L(w)\|}{\|\nabla\| \nabla L(w)\left\|_2\right\|_2}$
% % \STATE  \hspace{2cm} $g_{S_1,B_i} =\nabla^2 L(\theta+\delta^1) \cdot \frac{\nabla L(\theta+\delta^1)}{\|\nabla L(\theta+\delta^1)\|} = \nabla\| \nabla L(\theta+\delta^1)\|$
% \STATE  \hspace{2cm} $g_{S,B_i} = \nabla L_{S,B_i}(\theta)+ \nabla\| \nabla L(\theta+\delta^1)\|$
% \STATE \hspace{1.5cm} \textbf{FLAD}: 
% \STATE \hspace{2cm} \textbf{Accumulate} batch’s loss gradient to approximate full-batch gradient $g_{S} =  \lambda_0 g_{S} + (1-\lambda_0) g_{S,B}$
% \STATE \hspace{2cm} \textbf{ gradient not align with full-batch gradient} $\delta^0_{rand} = \rho \frac{g_{S,B}- \cos(g_{S},g_{S,B})g_{S}}{\|g_{S,B}- \cos(g_{S},g_{S,B})g_{S}\|+c}$
% \STATE \hspace{2cm} $g^0_{S,B_i}=\nabla L_{S,B_i}(\theta + \delta^0_{rand})$
% \STATE \hspace{2cm} \textbf{Accumulate} batch’s second gradient to approximate full-batch second gradient $\nabla \|g_S\| =  \lambda_1 \nabla \|g_S\| + (1-\lambda_1) \nabla \|g_{S,B}\|$

% \STATE \hspace{2cm} \textbf{second gradient not align with the full-batch direction} 
% $\delta^1_{rand}={\rho} \frac{\nabla\left\|{g}_{B_t}\right\|-\sigma \nabla \|g_S\|}{\|\nabla\| {g}_{B_t}\left\|-\sigma \nabla \|g_S\|\right\|+c}$
% \STATE \hspace{2cm} $g^1_{S,B_i}=\nabla\| \nabla L(\theta+\delta^1_{rand})\|$
% \STATE \hspace{2cm} $g_{S,B_i}=g_{S,B_i}^0 +g_{S,B_i}^1$

% % \STATE \hspace{1cm} \textbf{Using perturbation based mothod or SGD}
% % \STATE \hspace{1cm} \textbf{Accumulate} batch’s loss gradient to approximate full-batch gradient $g_{S} =  \lambda_0 g_{S} + (1-\lambda_0) g_{S_1,B}$
% % \STATE \hspace{1cm} \textbf{Compute zero-order random perturbation}(reduce align with gradient direction)$\delta_0 = \rho \frac{g_{S,B}- \cos(g_{S_1},g_{S_1,B})g_{S}}{\|g_{S,B}- \cos(g_{S_1},g_{S_1,B})g_{S}\|+c}$
% % \STATE \hspace{1cm} \textbf{Calculate gradient} from Loss with zero-order random perturation $g_{0}$
% \STATE \hspace{1cm} \textbf{Update model} $\theta_{S_1,B_{i+1}} = \theta_{S_1,B_{i}} - \eta  g_{S_1,B_i} $
% \STATE \hspace{1cm} {Add sample to sample memory buffer} $M_0 \leftarrow B$
% \STATE \hspace{0.5cm} \textbf{Merge Memory} $M \leftarrow M_0$  
% \STATE \hspace{0.5cm} \textbf{Save }  MODEL $\theta_1$, Memory $M$

% \STATE  \textbf{Other tasks} $t=2,\cdots,T$
% \STATE \hspace{0.5cm} \textbf{Batch} $i=1,\dots,N$

% \STATE \hspace{1cm}   \textbf{Batch} $i \%  K = 0$
% \STATE \hspace{1.5cm}   \textbf{Sample old task's sample} $S_{o}$ from buffer $M$
% \STATE \hspace{2cm}   \textbf{Calculate and Acumulate old gradient} $g_{old} = \lambda_0 g_{old} +(1-\lambda_0) g_{old,S}$
% \STATE \hspace{2cm}  \textbf{Calculate old second  gradient}  $\nabla \|g_{old}\| =  \lambda_1 \nabla \|g_{old}\| + (1-\lambda_1) \nabla \|g_{old,S_o}\|$

% \STATE \hspace{1.5cm} \textbf{Compute gradient} batch’s loss gradient $g_{S,B_i}=\nabla L_{S,B_i}(\theta)$


% \STATE \hspace{1.5cm} \textbf{Update model} $\theta_{S_t,B_{i+1}} = \theta_{S_t,B_{i}} - \eta  (g_{S_t,B_i} +g_{old})$

% \STATE \hspace{1cm}   \textbf{Batch} $i\% K \neq 0$



% \STATE \hspace{1.5cm} \textbf{Sample batch} Sample batch $B_i$ from $S_t$
% \STATE \hspace{1.5cm} \textbf{Compute gradient} batch’s loss gradient using optimizer and filiter with previous information
% \STATE \hspace{1.5cm} \textbf{SGD}: $g_{S,B_i}=\nabla L_{S,B_i}(\theta)$
% \STATE \hspace{2cm} $g_{S,B_i}=g_{S,B}- \cos(g_{old},g_{S,B})g_{S}$

% \STATE \hspace{1.5cm} \textbf{SAM} $g_{S,B_i}=\nabla L_{S,B_i}(\theta)$
 
% \STATE \hspace{2cm} \textbf{ gradient not align with previous gradient} $\delta^0_{new} = \rho \frac{g_{S,B}- \cos(g_{old},g_{S,B})g_{S}}{\|g_{S,B}- \cos(g_{old},g_{S,B})g_{S}\|+c}$
% \STATE \hspace{2cm} $g^0_{S,B_i}=\nabla L_{S,B_i}(\theta + \delta^0_{new})$
% \STATE \hspace{2cm} $g_{S,B_i}=g_{S,B}- \cos(g_{old},g_{S,B})g_{S}$

% \STATE \hspace{1.5cm} \textbf{FLAD-F}:
% \STATE \hspace{2cm} \textbf{ gradient not align with previous gradient} $\delta^0_{new} = \rho \frac{g_{S,B}- \cos(g_{old},g_{S,B})g_{S}}{\|g_{S,B}- \cos(g_{old},g_{S,B})g_{S}\|+c}$
% \STATE \hspace{2cm} $g^0_{S,B_i}=\nabla L_{S,B_i}(\theta + \delta^0_{new})$
% \STATE \hspace{2cm} \textbf{second gradient not align with previous second direction} 
% $\delta^1_{new}={\rho} \frac{\nabla\left\|{g}_{B_t}\right\|-\sigma \nabla \|g_S\|}{\|\nabla\| {g}_{B_t}\left\|-\sigma \nabla \|g_S\|\right\|+c}$
% \STATE \hspace{2cm} $g^1_{S,B_i}=\nabla\| \nabla L(\theta+\delta^1_{new})\|$
% \STATE \hspace{2cm} $g_{S,B_i}=g_{S,B_i}^0 +g_{S,B_i}^1$
% \STATE \hspace{2cm} $g'_{S,B_i}=g_{S,B}- \cos(g_{old},g_{S,B})g_{S,B_i}$


% \STATE \hspace{1.5cm} \textbf{Update model} $\theta_{S_t,B_{i+1}} = \theta_{S_t,B_{i}} - \eta  g_{S_t,B_i} $

% \STATE \hspace{1cm} {Add sample to sample memory buffer} $M_t \leftarrow B$
% \STATE \hspace{0.5cm} \textbf{Merge Memory} $M \leftarrow M_t$  

% \end{algorithmic}
% \label{alg1}
% \end{algorithm}

\section{Introduction}

% paragraph{Limitations of existing continual learning methods.}

% Most existing continual learning approaches regulate only the deterministic update direction.
% Regularization-based methods~\cite{EWC2017,onlineEWC,SI2017,KFLANEURIPS2018_Ritter}
% mitigate forgetting by penalizing parameter drift along directions that are important for previous tasks.
% These methods typically estimate parameter importance using curvature proxies such as the Hessian or Fisher information computed after each task.
% For instance, EWC~\cite{EWC2017} and Online-EWC~\cite{onlineEWC} employ the diagonal Fisher information matrix,
% while Synaptic Intelligence (SI)~\cite{SI2017} estimates importance from optimization trajectories.
% More accurate approximations such as Kronecker-factored Laplace approximation (KFLA)~\cite{KFLANEURIPS2018_Ritter} also incorporate off-diagonal curvature information.
% From an optimization perspective, these regularization mechanisms ultimately influence the gradient direction during training.
% A related line of work directly manipulates gradients during optimization.
% Gradient Episodic Memory (GEM)~\cite{lopez-pazGradientEpisodicMemory2017}
% and A-GEM~\cite{chaudhry2019efficientlifelonglearningagem}
% project the current gradient into a feasible region defined by gradients of episodic memories.
% Orthogonal Gradient Descent (OGD)~\cite{farajtabarOrthogonalGradientDescent2020}
% constructs a historical sensitivity subspace and removes its components from the update direction.
% Other optimization-based approaches similarly enforce constraints on gradient directions~\cite{NEURIPS2021_ec5aa0b7,NEURIPS2020_70d85f35,wang2023orthogonal}.
% Despite their differences, these methods share a common principle:
% they regulate \emph{drift} by modifying the gradient direction.

Most existing continual learning approaches regulate only the deterministic component of the optimization dynamics. In particular, both regularization-based and optimization-based methods ultimately control \emph{drift} by modifying the gradient direction during training. For example, regularization methods such as EWC and Online-EWC~\cite{EWC2017,onlineEWC} penalize parameter updates along directions that are important for previous tasks using Fisher information, while Synaptic Intelligence (SI)~\cite{SI2017} estimates similar importance from optimization trajectories. On the other hand, optimization-based methods directly manipulate gradients: GEM and A-GEM~\cite{lopez-pazGradientEpisodicMemory2017,chaudhry2019efficientlifelonglearningagem} project the current gradient into a feasible region defined by episodic gradients, while OGD~\cite{farajtabarOrthogonalGradientDescent2020} removes components that lie in a historical sensitivity subspace. Despite their different formulations, these approaches share a common principle: they regulate \emph{drift} by altering the gradient direction.

% \paragraph{Flatness-based continual learning.}

% Another line of work mitigates forgetting by promoting flat solutions in the loss landscape~\citep{NEURIPS2020_518a38cc,JMLRAnEmpiricalInvestigatio,liRevisitingCatastrophicForgetting2024}.
% Sharpness-aware optimization methods minimize neighborhood robust objectives~\cite{pmlr-v151-bisla22a,Zhang_GAM,NEURIPS2024_C-Flat,yangRevisitingFlatnessAwareOptimization2025,chen2026sharpnessflatnessdecompositionframework},
% while other approaches manipulate the loss landscape through weight averaging or mode connectivity~\cite{izmailov2019averagingweightsleadswider,mirzadeh2020linearmodeconnectivitymultitask}.
% Improved representations obtained from large-scale pretraining or self-supervised learning can also lead to flatter solutions~\cite{hu2022doesselfsupervisedpretrainingperform,ramasesh2022effect}.

% However, these methods typically encourage flatness with respect to the \emph{current task loss}.
% They do not explicitly consider how stochastic perturbations interact with the curvature of previous tasks.
% As a result, perturbations that help flatten the new-task loss may still move parameters toward directions that significantly increase the old-task loss.
% \paragraph{Flatness-based continual learning.}

Another line of work mitigates forgetting by promoting flat solutions in the loss landscape. These approaches typically encourage robustness of the \emph{current task loss} to local perturbations, but do not explicitly account for how stochastic perturbations interact with the curvature of previous tasks. 
Representative methods include sharpness-aware optimization techniques that minimize neighborhood-robust objectives~\cite{pmlr-v151-bisla22a,Zhang_GAM,NEURIPS2024_C-Flat,yangRevisitingFlatnessAwareOptimization2025,chen2026sharpnessflatnessdecompositionframework}, as well as approaches that manipulate the loss landscape through weight averaging or mode connectivity~\cite{izmailov2019averagingweightsleadswider,mirzadeh2020linearmodeconnectivitymultitask}. However, these methods primarily promote flatness with respect to the \emph{new task}. Consequently, perturbations that flatten the new-task loss may still align with high-curvature directions of previous tasks, potentially increasing their losses.


\section{Related Work}

\paragraph{Regularization and optimization CL methods }

Regularization-based methods~\cite{EWC2017,onlineEWC,SI2017,KFLANEURIPS2018_Ritter} mitigates forgetting by penalizing parameter drift along directions that are important for previous tasks,
% mitigates forgetting by adding an explicit consolidation term that penalizes parameter drift along directions that are deemed important for past tasks, 
% typical through the
% second-order Hessian of previous tasks computed at the end of each task training.
typically quantified using curvature proxies such as the Hessian or Fisher information computed after each task.
EWC~\cite{EWC2017} and Online-EWC~\cite{onlineEWC} employ the diagonal Fisher information matrix, 
% Intelligent Synapses (IS) 
Synaptic Intelligence (SI)~\cite{SI2017} estimates importance from within-task optimization trajectories,  Kronecker factored Laplace approximation
(KFLA)~\cite{KFLANEURIPS2018_Ritter} considering off-diagonal elements and utilizes a  Kronecker factorization approximation.
In essence, the constraint on continual learning for  regularization is ultimately reflected in gradient directions. 
Some recent work directly manipulates the gradient-based optimization process, categorized as the optimization-based approach~\cite{lopez-pazGradientEpisodicMemory2017,farajtabarOrthogonalGradientDescent2020,NEURIPS2021_ec5aa0b7,NEURIPS2020_70d85f35,wang2023orthogonal}. GEM~\cite{lopez-pazGradientEpisodicMemory2017} and A-GEM~\cite{chaudhry2019efficientlifelonglearningagem} projecting the current gradient into a feasible set defined by episodic gradients. OGD~\cite{farajtabarOrthogonalGradientDescent2020} instead constructs a historical sensitivity subspace and projects each update onto its orthogonal complement. NCL~\cite{NEURIPS2021_ec5aa0b7} manipulates gradient directions in online Laplace approximation.

% (Kirkpatrick et al., 2017; Ritter
% et al., 2018; Huszar´ , 2017; Zenke et al., 2017), being an important branch of these methods, aim to discern and retain the important parameter weights associated with prior tasks so as to sustain their
% performance. Specifically, they assess the importance of different parameter weights through the
% second-order Hessian of previous tasks computed at the end of each task training.


% The differences
% among these regularization-based methods primarily reside in their distinct approaches to approximating the Hessian. For example, EWC (Kirkpatrick et al., 2017) and Online-EWC (Huszar´ , 2017)
% employ the diagonal Fisher information matrix, whereas Kronecker factored Laplace approximation
% (KFLA) (Ritter et al., 2018) utilizes a more accurate Kronecker factorization approximation, considering off-diagonal elements. Similarly, Intelligent Synapses (IS) (Zenke et al., 2017) adopts an
% approximate diagonal matrix, but its elements take into account the importance of a task throughout
% the entire training trajectory. 


% Despite these methods striving to improve the Hessian approximation.
% In essence, the constraint on continual learning for either replay or regularization is ultimately reflected in gradient directions. As a result, some recent work directly manipulates the gradient-based optimization process, categorized as the optimization-based approach in Sec. 4.3. Alternatively, gradient projection can also be performed without storing old training samples [65], [115], [146], [202], [227], [258], [266], [333], [380], [448], [496]. he first term encourages parameter changes predominantly in directions that do not interfere with the old tasks via a preconditioner Λ−1  k−1, while the second term enforces θ to stay close to the o tasks via a preconditioner Λ−1  k−1, while the second term enforces θ to stay close to the old task solution μk−1..ld task solution μk−1.



% \paragraph{Flatness based CL methods}



\paragraph{Flatness based CL methods.}
Another line of work mitigates forgetting by explicitly promoting flat solutions~\citep{NEURIPS2020_518a38cc,JMLRAnEmpiricalInvestigatio,liRevisitingCatastrophicForgetting2024} through sharpness aware optimization that minimizes a neighborhood robust loss~\cite{pmlr-v151-bisla22a,Zhang_GAM,NEURIPS2024_C-Flat,yangRevisitingFlatnessAwareOptimization2025,chen2026sharpnessflatnessdecompositionframework}, loss landscape manipulation via ensembling or mode connectivity~\cite{izmailov2019averagingweightsleadswider,mirzadeh2020linearmodeconnectivitymultitask}, and improved representation learning via large scale pretraining or self supervised objectives~\cite{hu2022doesselfsupervisedpretrainingperform,ramasesh2022effect}. 
% Representative studies include sharpness aware CL (Bisla et al., 2022; Zhang et al., 2023; Bian et al., 2024; Yang et al., 2025; Chen et al., 2026), mode connectivity and averaging based approaches (Izmailov et al., 2019; Mirzadeh et al., 2020a), and representation centric methods (Hu et al., 2022; Ramasesh et al., 2022).
% It is showed improved flatness can mitigate forgetting in continual learning~\citep{NEURIPS2020_518a38cc,JMLRAnEmpiricalInvestigatio,liRevisitingCatastrophicForgetting2024}.
% However, the flatness related method actually only pay attention to the current work.  
% When the converged loss landscape is flatter and the observed data distributions are more similar, a properly balanced solution can better generalize to the task sequence . 


% To find such a flat minima, there are three widely-used strategies in continual learning. The first is derived from its definition, i.e., the flatness metric. Specifically, the minimization of lt(θ; Dt) can be replaced by a robust task-specific loss ltb(θ; Dt) := max∥δ∥≤b lt(θ + δ; Dt), thus the obtained solution guarantees low error not only at a specific point but also in its neighborhood with a “radius” of b. However, due to the great dimensionality of θ, the calculation of ltb(θ; Dt) cannot cover all possible δ but only a few directions [387], similar to the complexity issue of computing the Hessian matrix in Eq. (19). The alternatives include using an approximation of the Hessian [90], [313] or calculating δ only along the trajectory of forward and backward parameter changes [58], [174], [183], [300], [312]. The second is to operate the loss landscape by constructing an ensemble model under the restriction of mode connectivity, i.e., integrating multiple minima in parameter or function space along the low-error path, since connecting them ensures flatness on that path [60], [131], [183], [312], [443], [462]. These two strategies are closely related to the optimizationbased approach. The third comes down to obtaining welldistributed representations, which tend to be more robust to distribution differences in function space, such as by using pre-training [168], [300], [360], wider network architectures [309], [359], [360] and self-supervised learning [57], [168], [286], [343]. Observing the substantial attention to large-scale pre-training

% \paragraph{Optimization Dynamics.}
% A large body of work interprets the emergence of flat solutions under stochastic optimization via  stochastic dynamics process. Under suitable scaling, (stochastic) gradient methods can be approximated by an Stochastic Differential Equation (SDE) whose drift captures the deterministic descent direction and whose diffusion captures mini batch and algorithm-induced noise.
% Early studies suggest that even vanilla SGD can bias training toward relatively flat regions and generlization better due to the inherent noise in the gradient estimation.~\cite{NIPS2017_81b3833e_SGDgenerlizationbetter,keskar2017largebatchtrainingdeeplearning}. Zhu et al.~\cite{zhu2019anisotropicnoisestochasticgradient}analyze a general noisy gradient dynamics that unifies SGD and Langevin-type updates, showing that both escape efficiency and implicit regularization depend on the alignment between the noise covariance and the Hessian, which helps explain why anisotropic SGD noise can escape sharp minima more effectively. Sharpness-aware optimizers provide a more explicit mechanism for reshaping the dynamics. In particular, SAM can be viewed as modifying the drift using local curvature information, effectively biasing updates toward wider basins~\cite{JMLR:v24:23-043}. 

% However, most existing analyses do not connect these dynamics to forgetting. In continual learning, the central question is how \emph{old-task} risk evolves along \emph{new-task} training trajectories. Our process view leverages the same drift--diffusion decomposition, but applies it to an old-task energy, thereby isolating how optimization dynamics influence forgetting.

% % \paragraph{Optimization Dynamics.}
% % % A large body of work interprets stochastic optimization through continuous-time stochastic dynamics.
% % % Flatness is always a good property of the model but 
% % % the link between flatenss  and the traing process is unkown. 
% % A large body of work interprets  flatenss result and stochastic optimization process through continuous-time stochastic dynamics
% % Under suitable scaling, (stochastic) gradient methods can be approximated by an SDE whose drift captures the deterministic descent direction and whose diffusion captures minibatch and algorithm-induced noise. 


% % Early study show even SGD can produce a flat result.  Zhu et al.\ study a general noisy gradient dynamics unifying SGD and Langevin-type updates, and derive indicators showing that escape efficiency and implicit regularization depend on the alignment between the noise covariance and the Hessian, explaining why anisotropic SGD noise can escape sharp minima more effectively.
% % Sharpness-aware optimizers provide another explicit mechanism that reshapes optimization dynamics. 
% % In such cases, SAM’s
% % update may be regarded as a third derivative—the derivative of the Hessian in the leading
% % eigenvector direction—that encourages drift toward wider minima.

% % However, Most of the existing work igonre the relation with forgetting.
% % In contrast, continual learning requires tracking how \emph{old-task} risk evolves along \emph{new-task} training dynamics; our process view leverages the same drift-diffusion decomposition but applies  analysis to an old-task energy evaluated to figure out the influence of optimizaiton dynamics on forgetting.
% % % , enabling connections to stationary distributions and approximate Bayesian inference ~\citep{mandt2017sgd,chaudhari2018sgd,jastrzebski2018threefactors}. :contentReference[oaicite:0]{index=0}
% % % Beyond noise magnitude, recent analyses emphasize the \emph{structure} of the diffusion and its interaction with curvature.
% % % Zhu et al.\ study a general noisy gradient dynamics unifying SGD and Langevin-type updates, and derive indicators showing that escape efficiency and implicit regularization depend on the alignment between the noise covariance and the Hessian, explaining why anisotropic SGD noise can escape sharp minima more effectively than isotropic diffusion with matched variance ~\citep{zhu2019anisotropic}. 



% % % Sharpness-aware optimizers provide another explicit mechanism that reshapes optimization dynamics by evaluating gradients at locally perturbed points.
% % % Recent diffusion-based theories of SAM model its updates via an SDE and the associated Fokker--Planck equation, decomposing the probability current into drift and diffusion components; these works further connect SAM's gradient covariance matrix (GCM) to curvature through spectral ill-conditioning and strong GCM--Hessian alignment, and use Kramers-type arguments to relate this geometry to exponential escape behavior ~\citep{beyondsharpnessdensity}. 

% % % Related perturbation-based objectives such as gradient-norm aware minimization (GAM) formalize \emph{first-order flatness} via the maximal gradient norm in a local neighborhood and optimize it through multi-point gradient evaluations, again highlighting how local perturbations alter the effective drift and diffusion geometry ~\citep{gamfirstorderflatness}. :contentReference[oaicite:3]{index=3}


% % % Most of the above anal:contentReference[oaicite:4]{index=4}ingle} training energy (the task being optimized).











\paragraph{Optimization Dynamics.}
A large body of work explains the emergence of flat solutions under stochastic optimization through stochastic dynamics. Under suitable scaling, stochastic gradient methods admit an SDE approximation whose drift represents the deterministic descent direction and whose diffusion models minibatch and algorithm-induced noise. Early studies argue that this intrinsic noise can bias SGD toward flatter regions and improved generalization ~\cite{NIPS2017_81b3833e_SGDgenerlizationbetter,keskar2017largebatchtrainingdeeplearning}. Zhu et al.\ ~\cite{zhu2019anisotropicnoisestochasticgradient} analyze a unified noisy-gradient dynamics covering SGD and Langevin-type updates, and show that escape efficiency and implicit regularization depend on the alignment between the noise covariance and the Hessian, which clarifies why anisotropic noise can escape sharp minima more effectively. Sharpness-aware optimizers provide a more explicit mechanism: SAM can be interpreted as a curvature-informed modification of the drift that biases updates toward wider basins ~\cite{JMLR:v24:23-043}.

Most existing analyses, however, do not address forgetting. In continual learning, the key quantity is how \emph{old-task} risk evolves along \emph{new-task} training dynamics. Our process view uses the same drift--diffusion decomposition, but applies it to an old-task energy, thereby isolating the influence of optimization dynamics on forgetting.





\section{Optimization Dynamics of Forgetting}

In continual learning, catastrophic forgetting arises because the optimization dynamics of a new task modify parameters in directions that increase the loss of previous tasks.
Understanding this phenomenon requires analyzing the interaction between optimization trajectories and training-time stochasticity.

This section develops a process-level view of forgetting.
We first analyze the training trajectory using stochastic differential equations and derive a decomposition of the change in the old-task risk.
We then connect this process view to a distributional perspective based on parameter perturbations and flatness.
These two views reveal a common quantity governing forgetting: the curvature–covariance interaction term
$tr(H\Sigma)$.
Finally, this observation leads to a principled control objective for continual learning.




% In continual learning tasks arrive sequentially.
% Let $\mathcal{X}$ denote the input space, $\mathcal{Y}$ the label space, and $\mathcal{Z}=\mathcal{X}\times\mathcal{Y}$.
% A model with parameters $W\in\mathbb{R}^{d}$ induces a predictor $f_W:\mathcal{X}\rightarrow\mathcal{Y}$.



% For each task $t$, the CL system (which we will call an agent) receives a training set $S_t = \{(x_{tj}, y_{tj})\}_{j=1}^{m_t}$ and a test set $T_t$, both drawn i.i.d.\ from a task-specific data distribution $\mathcal{D}_t$. The empirical loss for function $f \in \mathcal{F}$ on task $t$ is defined as
% $
% \hat{L}_{S_t}(f) = \frac{1}{m_t} \sum_{j=1}^{m_t} \ell(f(x_{tj}), y_{tj}).
% $

% Given a dataset $S={z_j}_{j=1}^{m}$ with $z_j\in\mathcal{Z}$, the empirical risk is $\hat L_S(W)=\frac{1}{m}\sum_{j=1}^{m}\ell(W,z_j)$. 
% Let $\theta_{t-1}$ denote the parameter after finishing task $t-1$ and $\theta_t$ the parameter obtained after training on task $t$.

% % \begin{equation}
% % \hat L_S(W)=\frac{1}{m}\sum_{j=1}^{m}\ell(W,z_j).
% % \end{equation}




% We analyze forgetting on task $1$ induced by the training dynamics of task $2$.
% Let the parameter trajectory during task-2 training be ${\theta}_{t\in[0,T]}$ with $\theta_0=\theta_1, \theta_T=\theta_2 .$
% % \begin{equation}
% % \theta_0=\theta_1,
% % \qquad
% % \theta_T=\theta_2 .
% % \end{equation}
% The parameter displacement is $\delta_{1\rightarrow2}=\theta_2-\theta_1 .
% $
% % \begin{equation}
% % \delta_{1\rightarrow2}=\theta_2-\theta_1 .
% % \end{equation}
% To analyze forgetting during task-2 training,  we track the old-task loss $\hat L_{S_1}(\theta)$ along this trajectory and define forgetting is 
% \begin{equation}
%     \mathbb{E}[L_{S_1}(\theta_2)]-L_{S_1}(\theta_1).
% \end{equation}
\subsection{Forgetting in the dynamic training process}

In continual learning, tasks arrive sequentially.
Let $\mathcal{X}$ denote the input space, $\mathcal{Y}$ the label space, and $\mathcal{Z}=\mathcal{X}\times\mathcal{Y}$ the sample space.
A model with parameters $W\in\mathbb{R}^d$ induces a predictor $f_W:\mathcal{X}\rightarrow\mathcal{Y}$.
At task $t$, the learner receives a training dataset
$
S_t={(x_{tj},y_{tj})}_{j=1}^{m_t}
$
and a test set $T_t$, both drawn i.i.d. from a task specific distribution $\mathcal{D}_t$ over $\mathcal{Z}$.
For a dataset $S={(x_j,y_j)}_{j=1}^{m}$, the empirical risk of parameter $W$ is $\hat L_S(W)=
\frac{1}{m}
\sum_{j=1}^{m}
\ell\big(f_W(x_j),y_j\big).$
% \begin{equation}
% \hat L_S(W)=
% \frac{1}{m}
% \sum_{j=1}^{m}
% \ell\big(f_W(x_j),y_j\big).
% \end{equation}

Let $\theta_{t-1}$ denote the model parameters obtained after completing task $t-1$.
Training on task $t$ updates the parameters through a learning algorithm $\mathcal{A}_t$, $\theta_t = \mathcal{A}_t(\theta_{t-1}, S_t)$.
% \begin{equation}
% \theta_t = \mathcal{A}t(\theta{t-1}, S_t).
% \end{equation}
% This sequential update typically leads to catastrophic forgetting, where learning a new task degrades performance on previous ones.
Such sequential updates often lead to catastrophic forgetting, where learning a new task deteriorates performance on previously learned tasks.
In this work we analyze forgetting on task $1$ induced by the training dynamics of task $2$.
Let $\theta_1$ denote the parameter obtained after completing task $1$, and $\theta_2$ the parameter obtained after training on task $2$.
We quantify forgetting as the decrease of the old-task empirical loss.
% Let $\theta_1$ denote the parameter obtained after finishing task $1$, and $\theta_2$ the parameter after training on task $2$.
% We define forgetting as the Loss decrease of the previous task  to measures the degradation of the task-1 objective caused by subsequent training on task 2
\begin{definition}(Forgetting)
The forgetting from task $1$ to task $2$ is defined as
    \begin{equation}
\Delta_{1\rightarrow2}=
\hat L_{S_1}(\theta_2)
-
\hat L_{S_1}(\theta_1).
\end{equation}
\end{definition}
% This quantity measures the change of the task-1 objective after training on task 2.
% In practice, the value is typically negative, indicating performance degradation on task $1$ induced by learning task $2$.
This notion is closely related to the classical Backward Transfer (BWT) metric in continual learning. The main difference is that BWT is defined on the $0$–$1$ evaluation loss (or accuracy) measured on test data, whereas our formulation is based on the empirical training loss $\hat L_{S_1}(\theta)$. 
% This choice allows us to analyze forgetting through the lens of optimization dynamics and stochastic training processes.
% This quantity measures the degradation of the task-1 objective caused by subsequent training on task 2.
%track the old task loss $\hat L_{S_1}(W)$ along the optimization trajectory and


% A standard way to minimize the empirical risk $\hat L_{S_2}(\theta)$ is 
% gradient based optimizers act like GD with an unbiased noise, including gradient Langevin dynamics (GLD) and and stochastic gradient descent (SGD).
% The update can be written as the dynamics of gradient descent with unbiased noise~\cite{zhu2019anisotropicnoisestochasticgradient}:
A standard approach to minimize the empirical risk $\hat L_{S_2}(\theta)$ is gradient based optimization.
Many practical optimizers, including stochastic gradient descent (SGD) and gradient Langevin dynamics (GLD), can be modeled as gradient descent perturbed by unbiased noise ~\cite{zhu2019anisotropicnoisestochasticgradient}.The parameter update can therefore be written as
\begin{equation}
\theta_{t+1}
=
\theta_t
-
\eta\,\nabla_{\theta} \hat L_S(\theta_t)
+
\eta\,\epsilon_t,
\quad \epsilon_t\sim \mathcal{N}(0,\Sigma(\theta_t))
% \qquad
% \mathbb E[\epsilon_k\mid \theta_k]=0,
% \qquad
% \operatorname{Cov}(\epsilon_k\mid \theta_k)=\Sigma_k .
\label{eq:noisy_gd_discrete_explain}
\end{equation}
% Here $\eta>0$ is the step size (learning rate), $\Sigma_t$ is distribution of noise in the optimal process. 
where $\eta>0$ denotes the step size and $\Sigma(\theta_t)$ represents the covariance of the optimization noise.
% and $\Sigma^{sgd}(\theta_t) \approx \frac{1}{m}\left[\frac{1}{N} \sum_{i=1}^N \nabla \ell\left(x_i ; \theta_t\right) \nabla \ell\left(x_i ; \theta_t\right)^T-\nabla L\left(\theta_t\right) \nabla L\left(\theta_t\right)^T\right]$.
% For small enough constant learning rate $\eta$, Eq. \eqref{eq:noisy_gd_discrete_explain} can be treated as the numerical discretization of the following stochastic differential equation
For sufficiently small learning rate $\eta$, Eq.~\eqref{eq:noisy_gd_discrete_explain} can be viewed as a numerical discretization of the stochastic differential equation(SDE)
\begin{equation}
d\theta_t
=
- q_t\,dt
+
\Sigma_t^{1/2}\,dW_t,
\label{eq:noisy_gd_sde_explain}
\end{equation}
% where $W_t$ is a standard Brownian motion,$\Sigma_t$ is the effective diffusion capturing the covariance of training-time noise. $q_t$ denotes the optimizer-shaped descent direction on task 2. In the simplest case, $q_t=\nabla \hat L_{S_2}(\theta_t)$. 
where $W_t$ is a standard Brownian motion and $\Sigma_t^{1/2} = \sqrt{\eta\Sigma(\theta_t)}$ denotes the effective diffusion capturing the covariance of training-time noise.
The vector $q_t$ represents the optimizer-induced descent direction for task $2$. In the simplest case, $q_t=\nabla_{\theta} \hat L_{S_2}(\theta_t)$.
% Provided that the mild smoothness assumptions for Ito’s lemma ~\cite{oksendal2003stochastic} holds, we have the optimization process of the train on task $2$
Under standard smoothness conditions required by Itô's lemma ~\cite{oksendal2003stochastic}, the optimization dynamics for task $2$ satisfy
~\cite{zhu2019anisotropicnoisestochasticgradient}
\begin{equation}
    \mathbb{E}\left[L_{S_2}\left(\theta_t\right)-L_{S_2}\left(\theta_0\right)\right]=-
    \underbrace{\int_0^t \mathbb{E}\left[q_t^T q_t\right]}_{\text{drift}}
    +\underbrace{\int_0^t \frac{1}{2} \mathbb{E} \operatorname{tr}\left(H_{2,t} \Sigma^{\mathrm{alg}}_{S_2}(\theta_t)\right) \mathrm{d}t}_{\text{diffusion}},
\label{eq: current task energy}
\end{equation}
where $H_{2,t}=\nabla^2 \hat L_{S_2}(\theta_t)$ denotes the Hessian of the task-2 loss.




% $H_{2,t}:=\nabla^2 \hat L_{S_2}(\theta_t)$ is the Hessian of $L_{S_2}(\theta_t)$. 

% Stochastic gradient methods can be modeled by a stochastic differential equation

% \begin{equation}
% d\theta_t=-q_t,dt+\Sigma_t^{1/2}dW_t ,
% \end{equation}

% where $q_t$ denotes the optimizer-shaped descent direction on the new task and $\Sigma_t$ represents the covariance of training-time noise.


% Applying Itô’s lemma yields
% \begin{equation}
% \frac{d}{dt}\mathbb{E}[L_1(\theta_t)]
% =
% -\mathbb{E}\left[\nabla L_1(\theta_t)^\top q_t\right]
% +
% \frac12\mathbb{E}\left[tr(H_{1,t}\Sigma_t)\right],
% \end{equation}

The result above characterizes the optimization dynamics of the current task objective $\hat L_{S_2}(\theta)$ under noisy gradient updates.
However, continual learning concerns not only the improvement of the current task, but also the effect of this optimization process on previously learned tasks. To analyze forgetting, we examine how the old-task loss $\hat L_{S_1}(\theta)$ evolves along the same parameter trajectory ${\theta_t}{t\in[0,T]}$ generated by training on task 2.
% In other words, although the optimization is driven by the gradient of $\hat L_{S_2}$, the parameters simultaneously move on the loss landscape of task 1.
The change of $\hat L_{S_1}(\theta_t)$ along this trajectory captures the degradation of the previous task induced by learning the new task.

Applying Itô's lemma to $\hat L_{S_1}(\theta_t)$ under the stochastic dynamics in Eq.~\eqref{eq:noisy_gd_sde_explain} yields the following pathwise decomposition of forgetting.



\begin{theorem}
    Assume $L_{S_1}(\theta)$ is twice continuously differentiable and the parameter trajectory follows the stochastic dynamics in Eq.~\eqref{eq:noisy_gd_sde_explain}. Then the expected forgetting from task $1$ to task $2$ satisfies
    \begin{equation}
    \mathbb{E}[\Delta_{1\rightarrow 2}]=
\mathbb{E}[L_{S_1}(\theta_t)-L_{S_1}(\theta_0)]
=
- \underbrace{\int_0^T
\mathbb{E}[v_t^\top q_t]dt}_{\text{drift interference}}
+
\underbrace{\frac12
\int_0^T
\mathbb{E}[tr\left(H_{1,t}\Sigma^{alg}_{S_2}(\theta_t)\right)]dt}_{\text{diffusion interference}} .
\label{eq: current task forgetting}
\end{equation} 
where $v_t=\nabla L_{S_1}(\theta_t),\quad H_{1,t}=\nabla^2 L_{S_1}(\theta_t)$. 
% The first term corresponds to gradient interference, while the second term captures curvature-weighted noise during training.
\end{theorem}

% Integrating over the training trajectory gives the change in old-task loss
% \begin{equation}
% \mathbb{E}[L_{S_1}(\theta_2)-L_{S_1}(\theta_1)]
% =
% - \underbrace{\int_0^T
% \mathbb{E}[\nabla L_{S_1}(\theta_t)^\top q_t]dt}_{\text{drift}}
% +
% \underbrace{\frac12
% \int_0^T
% \mathbb{E}[tr(H_{1,t}\Sigma_t)]dt}_{\text{diffusion}} .
% \end{equation} 
% where $H_{1,t}=\nabla^2 L_{S_1}(\theta_t)$.
% This decomposition reveals two mechanisms governing forgetting:
The decomposition reveals two fundamental mechanisms governing forgetting during optimization:
\begin{itemize}
    \item  Drift interference $\nabla L_1(\theta_t)^\top q_t ,$ which measures how the deterministic update direction of task $2$ affects the objective of task $1$.
    \item Diffusion interference $tr(H_{1,t}\Sigma_t)$ which quantifies how training-time noise interacts with the curvature of the old-task loss.
\end{itemize} 
Together, these two terms provide a dynamical perspective on catastrophic forgetting: the first captures gradient interference across tasks, while the second characterizes curvature weighted noise amplification along the training trajectory.
% \begin{equation}
% \nabla L_1(\theta_t)^\top q_t ,
% \end{equation}
% **Diffusion interference**
% \begin{equation}
% tr(H_{1,t}\Sigma_t) ,
% \end{equation}
% which measures how training-time noise interacts with the curvature of the old-task loss.

% This process-level decomposition provides a dynamical explanation of forgetting during optimization.
% This decomposition reveals two mechanisms governing forgetting:

% Equation \eqref{eq:path_integral_oldtask} exposes two distinct channels governing old task degradation:
% \begin{itemize}
% \item \textbf{Drift interference.}
% The integral of $\nabla \hat L_{S_1}(\theta_t)^{\top}q_t$ quantifies how the task 2 descent direction changes the old task loss along the trajectory.
% \item \textbf{Diffusion interference.}
% The integral of $\operatorname{tr}\big(H_{1,t}\Sigma^{\mathrm{alg}}_{S_2}(\theta_t)\big)$ quantifies how training time diffusion injects variance along old task sensitive directions.
% \end{itemize}


\subsection{Forgetting in Distributional Perspective}


During stochastic optimization, minibatch sampling and algorithmic randomness introduce perturbations to the parameters. 
These perturbations explore a local neighborhood around the current parameter $\theta$. 
To quantify the sensitivity of the loss landscape to such perturbations, we introduce randomized sharpness.

\begin{definition}[Randomized Sharpness]
For Gaussian perturbation $\varepsilon \sim \mathcal N(0,\Sigma)$,
\begin{equation}
\mathrm{RS}_S^{(0)}(\theta,\Sigma)
=
\mathbb E_{\varepsilon\sim\mathcal N(0,\Sigma)}
\big[
\hat L_S(\theta+\varepsilon)
\big]
-
\hat L_S(\theta).
\end{equation}
\end{definition}

Randomized sharpness characterizes the local exploration induced by the diffusion noise of stochastic optimization. 
When the noise covariance aligns with the curvature of the loss landscape, the stochastic dynamics escape sharp directions more efficiently and converge toward flatter regions ~\cite{zhu2019anisotropicnoisestochasticgradient}.

Under the stochastic dynamics in Eq.~\eqref{eq:noisy_gd_sde_explain}, the parameter trajectory $\theta_t$ forms a stochastic process that induces a time-dependent parameter distribution $P_t$. 
This observation motivates a distributional view of optimization.

For any parameter distribution $P$, we define the distributional empirical risk
\[
\mathcal L_S(P)
=
\mathbb E_{W\sim P}[\hat L_S(W)] .
\]
In practice, the local parameter distribution induced by stochastic optimization can be approximated by a Gaussian $P=\mathcal N(\theta,\Sigma)$.
% \[
% P=\mathcal N(\theta,\Sigma).
% \]

\begin{proposition}[Distributional Sharpness]
Assume $\hat L_S(\theta)$ is twice continuously differentiable. 
For $P=\mathcal N(\theta,\Sigma)$,
\[
\mathcal L_S(P)
\approx
\hat L_S(\theta)
+
\frac12 \operatorname{tr}(H_S(\theta)\Sigma),
\]
where $H_S(\theta)=\nabla^2 \hat L_S(\theta)$. 
Consequently,
\[
\mathrm{RS}_S^{(0)}(\theta,\Sigma)
\approx
\frac12 \operatorname{tr}(H_S(\theta)\Sigma).
\]
\end{proposition}

This result establishes a connection between randomized sharpness and the distributional risk induced by stochastic optimization. 
We now use this perspective to analyze forgetting from a distributional viewpoint.














% The previous analysis characterizes forgetting along a stochastic optimization trajectory using the It\^{o} decomposition. 
% We now provide a complementary view based on parameter perturbations and the induced parameter distribution.

% During training, stochastic optimization introduces random perturbations to the parameters due to minibatch sampling and algorithmic noise. 
% These perturbations can be interpreted as exploring a local neighborhood around the current parameter $\theta$. 
% To quantify the sensitivity of the loss landscape to such perturbations, we introduce the notion of random sharpness.

% \begin{definition}[Random Sharpness]
% For Gaussian perturbation $\varepsilon \sim \mathcal N(0,\Sigma)$, the random sharpness is defined as
% \begin{equation}
% \mathrm{RS}_S^{(0)}(\theta,\Sigma)
% =
% \mathbb E_{\varepsilon\sim\mathcal N(0,\Sigma)}
% \big[
% \hat L_S(\theta+\varepsilon)
% \big]
% -
% \hat L_S(\theta).
% \end{equation}
% \end{definition}

% This quantity measures the sensitivity of the empirical loss to perturbations within a local neighborhood of $\theta$ and therefore characterizes the local geometry of the loss landscape.

% From a stochastic optimization perspective, the parameter trajectory $\theta_t$ follows the stochastic dynamics in Eq.~\eqref{eq:noisy_gd_sde_explain}. 
% Consequently, $\theta_t$ forms a stochastic process that induces a time-dependent parameter distribution $p_t(\theta)$ governed by the corresponding Fokker--Planck equation ~\cite{Risken1996,TheVariationalFormulationoftheFokkerEquation}. 
% Instead of analyzing a single trajectory, we  therefore study the statistics of the induced parameter distribution.
% % \begin{definition}
% For any parameter distribution $P$, we define the distributional empirical risk on dataset $S$ as
% $
% \mathcal L_S(P)
% =
% \mathbb E_{W\sim P}[\hat L_S(W)] .
% $
% % \end{definition}
% In practice, the local distribution induced by stochastic optimization can be approximated by a Gaussian distribution
% $P = \mathcal N(\theta,\Sigma),$
% where $\theta$ denotes the mean parameter and $\Sigma$ the covariance.
% \begin{proposition}[Distributional Sharpness]
% Assume $\hat L_S(\theta)$ is twice continuously differentiable. 
% For a Gaussian parameter distribution $P=\mathcal N(\theta,\Sigma)$, the distributional risk admits the local approximation
% $
% \mathcal L_S(P)
% \approx
% \hat L_S(\theta)
% +
% \frac12 \operatorname{tr}\!\big(H_S(\theta)\Sigma\big),
% $
% where $H_S(\theta)=\nabla^2 \hat L_S(\theta)$.
% Consequently, random sharpness admits the approximation
% \begin{equation}
% \mathrm{RS}_S^{(0)}(\theta,\Sigma)
% \approx
% \frac12 \operatorname{tr}\!\big(H_S(\theta)\Sigma\big).
% \end{equation}
% \end{proposition}
% Random sharpness therefore characterizes the local exploration induced by the diffusion term of stochastic optimization. 
% When the noise covariance is aligned with the curvature of the loss landscape, the resulting stochastic dynamics escape sharp directions more efficiently and converge to flatter regions, consistent with the anisotropic noise analysis of SGD ~\cite{zhu2019anisotropicnoisestochasticgradient}. 












% -----------------------

% The previous analysis characterizes forgetting along a stochastic optimization trajectory using the It\^{o} decomposition.  
% We now present an equivalent formulation from a distributional perspective. 
% Instead of tracking a single trajectory $\theta_t$, we analyze the evolution of the parameter distribution induced by stochastic optimization.

% Under the stochastic dynamics in Eq.~\eqref{eq:noisy_gd_sde_explain}, the parameter sequence $\theta_t$ forms a stochastic process that induces a time-dependent distribution $p_t(\theta)$. 
% This distribution evolves according to the Fokker--Planck equation ~\cite{Risken1996,TheVariationalFormulationoftheFokkerEquation} associated with the stochastic differential equation.

% For any parameter distribution $P$, we define the distributional empirical risk on dataset $S$ as
% \[
% \mathcal L_S(P)
% =
% \mathbb E_{W\sim P}[\hat L_S(W)] .
% \]

% In practice, the local parameter distribution induced by stochastic optimization can be approximated by a Gaussian distribution
% \[
% P = \mathcal N(\theta,\Sigma),
% \]
% where $\theta$ denotes the mean parameter and $\Sigma$ the covariance.

% \begin{definition}[Random Sharpness]
% For Gaussian perturbation $\varepsilon \sim \mathcal N(0,\Sigma)$, the random sharpness is defined as
% \begin{equation}
% \mathrm{RS}_S^{(0)}(\theta,\Sigma)
% =
% \mathbb E_{\varepsilon\sim\mathcal N(0,\Sigma)}
% \big[
% \hat L_S(\theta+\varepsilon)
% \big]
% -
% \hat L_S(\theta).
% \end{equation}
% \end{definition}

% \begin{proposition}[Distribution Sharpness]
% Assume $\hat L_S(\theta)$ is twice continuously differentiable. 
% For a Gaussian parameter distribution $P=\mathcal N(\theta,\Sigma)$, the distributional risk admits the local approximation
% \begin{equation}
% \mathcal L_S(P)
% \approx
% \hat L_S(\theta)
% +
% \frac12 \operatorname{tr}\!\big(H_S(\theta)\Sigma\big),
% \end{equation}
% where $H_S(\theta)=\nabla^2 \hat L_S(\theta)$.
% Consequently,
% \begin{equation}
% \mathrm{RS}_S^{(0)}(\theta,\Sigma)
% \approx
% \frac12 \operatorname{tr}\!\big(H_S(\theta)\Sigma\big).
% \end{equation}
% \end{proposition}

\begin{definition}[Distributional Forgetting]
Let $P_1=\mathcal N(\theta_1,\Sigma_1)$ and $P_2=\mathcal N(\theta_2,\Sigma_2)$ denote the parameter distributions after tasks 1 and 2. 
The distributional forgetting on the old task $S_1$ is defined as
\begin{equation}
\Delta_{1\rightarrow2}^{\mathrm{dist}}
=
\mathcal L_{S_1}(P_2)
-
\mathcal L_{S_1}(P_1).
\end{equation}
\end{definition}

Using Proposition 1, distributional forgetting admits the decomposition
\begin{equation}
\Delta_{1\rightarrow2}^{\mathrm{dist}}
=
\underbrace{
\hat L_{S_1}(\theta_2)-\hat L_{S_1}(\theta_1)
}_{\text{mean drift}}
+
\underbrace{
\mathrm{RS}_{S_1}^{(0)}(\theta_2,\Sigma_2)
-
\mathrm{RS}_{S_1}^{(0)}(\theta_1,\Sigma_1)
}_{\text{sharpness change}} .
\end{equation}

\begin{theorem}[Path--Distribution Equivalence]
Assume $\hat L_{S_1}(\theta)$ is twice continuously differentiable and the parameter trajectory follows the stochastic dynamics in Eq.~\eqref{eq:noisy_gd_sde_explain}.  
Let $P_t$ denote the distribution of $\theta_t$.
Then the expected forgetting derived in Theorem~1 admits an equivalent distributional representation
\begin{equation}
\mathbb E[\Delta_{1\rightarrow2}]
=
\mathcal L_{S_1}(P_T)
-
\mathcal L_{S_1}(P_0),
\end{equation}
where $\mathcal L_S(P)=\mathbb E_{W\sim P}[\hat L_S(W)]$ denotes the distributional empirical risk.

Therefore, the trajectory-based It\^{o} decomposition characterizes the same quantity as the change of the distributional risk induced by the evolution of the parameter distribution.
\end{theorem}











% \begin{theorem}[Path--Distribution Equivalence]
% Assume the parameter trajectory follows the stochastic dynamics in Eq.~\eqref{eq:noisy_gd_sde_explain}. 
% Let $P_t$ denote the distribution of $\theta_t$. 
% Then the expected forgetting on the old task satisfies
% \begin{equation}
% \mathbb E[\Delta_{1\rightarrow2}]
% =
% -
% \int_0^T
% \mathbb E[
% \nabla \hat L_{S_1}(\theta_t)^\top q_t
% ] dt
% +
% \frac12
% \int_0^T
% \mathbb E[
% \operatorname{tr}(H_{1,t}\Sigma_t)
% ] dt .
% \end{equation}

% Moreover, the same quantity can be expressed as the change in distributional risk
% \begin{equation}
% \mathbb E[\Delta_{1\rightarrow2}]
% =
% \mathcal L_{S_1}(P_T)
% -
% \mathcal L_{S_1}(P_0).
% \end{equation}
% \end{theorem}

% \begin{corollary}[Flatness--Forgetting Relation]
% Under the quadratic approximation,
% \begin{equation}
% \operatorname{tr}(H_{1,t}\Sigma_t)
% =
% 2\,\mathrm{RS}_{S_1}^{(0)}(\theta_t,\Sigma_t).
% \end{equation}

% Therefore the diffusion component of forgetting can be written as
% \begin{equation}
% \frac12
% \int_0^T
% \operatorname{tr}(H_{1,t}\Sigma_t) dt
% =
% \int_0^T
% \mathrm{RS}_{S_1}^{(0)}(\theta_t,\Sigma_t) dt .
% \end{equation}
% \end{corollary}

% The above results show that expected sharpness characterizes the instantaneous diffusion-induced forgetting along the optimization trajectory. 
% Consequently, flat solutions with small expected sharpness reduce the accumulation of forgetting during training.





% ---

% \subsection{Forgetting in Distributional Perspective}

% % The previous analysis characterizes forgetting along a single optimization trajectory.
% % To analyze sensitivity to parameter perturbations, we adopt a distributional view.
% The previous analysis characterizes forgetting along a stochastic optimization trajectory using the Itô decomposition.
% We now present an equivalent distributional formulation that analyzes the same phenomenon through the statistics of the parameter distribution induced by stochastic optimization

% Under the stochastic dynamics(SDE) in Eq.\eqref{eq:noisy_gd_sde_explain}, the parameter sequence $\theta_t$ forms a stochastic process and induces a time dependent distribution $p_t(\theta)$.
% This distribution evolves according to the Fokker–Planck equation~\cite{Risken1996,TheVariationalFormulationoftheFokkerEquation} associated with the above stochastic differential equation.


% For any parameter distribution $P$, we define the distributional empirical risk on dataset $S$ as $\mathcal L_S(P)
% =
% \mathbb E_{W\sim P}\big[\hat L_S(W)\big].$
% % \begin{equation}
% % \mathcal L_S(P)
% % =
% % \mathbb E_{W\sim P}\big[\hat L_S(W)\big].
% % \end{equation}
% In practice the local parameter distribution induced by stochastic optimization can be approximated by a Gaussian distribution
% $P = \mathcal N(\theta,\Sigma)$
% % \begin{equation}
% % P = \mathcal N(\theta,\Sigma),
% % \end{equation}
% where $\theta$ denotes the mean parameter and $\Sigma$ the covariance of the perturbation.
% % \begin{Proposition}[Distributional Sharpness]  
% % \end{Proposition}
% % Let $P=\mathcal N(\theta,\Sigma)$ be a local Gaussian parameter distribution.

% \begin{definition}[randomized sharpness]
% % For Gaussian perturbation $\varepsilon\sim\mathcal N(0,\Sigma)$ 
% \begin{equation}
% \mathrm{RS}_S^{(0)}(\theta,\Sigma)=
% \mathbb E_{\varepsilon\sim\mathcal N(0,\Sigma)}
% \big[
% \hat L_S(\theta+\varepsilon)
% \big]-
% \hat L_S(\theta) 
% \end{equation}
% \end{definition}
% % we define randomized sharpness
% % \begin{equation}
% % \mathrm{RS}_S^{(0)}(\theta,\Sigma)=
% % \mathbb E_{\varepsilon}
% % \big[
% % \hat L_S(\theta+\varepsilon)
% % \big]-
% % \hat L_S(\theta)
% % \end{equation}
% % which 

% Under the quadratic approximation reduces to

% \begin{lemma}
% \begin{equation}
% \mathrm{RS}_S^{(0)}(\theta,\Sigma)
% \approx
% \frac12
% \operatorname{tr}\big(H_S(\theta)\Sigma\big).
% \end{equation}
% \end{lemma}
% \begin{proof}
%     Assume $\hat L_S(\theta)$ is twice continuously differentiable.
% Then the distributional risk admits the second order approximation

% \begin{equation}
% \mathcal L_S(P)
% \approx
% \hat L_S(\theta)
% +
% \frac12
% \operatorname{tr}\big(H_S(\theta)\Sigma\big),
% \end{equation}

% where $H_S(\theta)=\nabla^2\hat L_S(\theta)$.

% The additional term

% \begin{equation}
% \frac12 \operatorname{tr}(H_S(\theta)\Sigma)
% \end{equation}

% quantifies the sensitivity of the loss to parameter perturbations and corresponds to expected sharpness.
% \end{proof}


% % to the optimization states after tasks 1 and 2.
% We define  forgetting in the distribution view as
% \begin{definition}[Distributional forgetting]
% Let $P_1=\mathcal{N}(\theta_1,\Sigma_1)$ and $P_2=\mathcal{N}(\theta_2,\Sigma_2)$ to the optimization states after tasks 1 and 2. The distributional forgetting on the old task $S_1$ is
% \begin{equation}
% \Delta_{1\rightarrow2}^{\mathrm{dist}}=
% \mathcal L_{S_1}(P_2)
% -
% \mathcal L_{S_1}(P_1).
% \end{equation}
% \end{definition}


% Using Proposition 1 yields

% \begin{equation}
% \begin{aligned}
% \Delta_{1\rightarrow2}^{\mathrm{dist}}
% &=
% \underbrace{
% \hat L_{S_1}(\theta_2)-
% \hat L_{S_1}(\theta_1)
% }_{\text{mean drift}}
% +
% \underbrace{
% \mathrm{RS}^{(0)}_{S_1}(\theta_2,\Sigma_2)-
% \mathrm{RS}^{(0)}_{S_1}(\theta_1,\Sigma_1)
% }_{\text{sharpness change}} .
% \end{aligned}
% \end{equation}

% The first term captures the change in the mean parameter, while the second term reflects how the covariance of the parameter distribution interacts with the curvature of the old task.

% \begin{theorem}[Path–Distribution Equivalence]
%     \begin{equation}
% \mathbb E[\Delta_{1\rightarrow2}]
% =
% \int_0^T
% \mathbb E
% \big[
% \nabla \hat L_{S_1}(\theta_t)^\top q_t
% \big] dt
% +
% \frac12
% \int_0^T
% \mathbb E
% \big[
% \operatorname{tr}(H_{1,t}\Sigma_t)
% \big] dt .
% \end{equation}

% Moreover, this quantity equals the change of the distributional risk

% \begin{equation}
% \mathbb E[\Delta_{1\rightarrow2}]
% \mathcal L_{S_1}(P_T)

% \mathcal L_{S_1}(P_0),
% \end{equation}
% \end{theorem}
% % Theorem 2 (Path–Distribution Equivalence)

% Assume the parameter trajectory follows the stochastic dynamics.Let $p_t(\theta)$ denote the induced parameter distribution.
% Then the expected forgetting on the old task satisfies

% \begin{equation}
% \mathbb E[\Delta_{1\rightarrow2}]
% =
% \int_0^T
% \mathbb E
% \big[
% \nabla \hat L_{S_1}(\theta_t)^\top q_t
% \big] dt
% +
% \frac12
% \int_0^T
% \mathbb E
% \big[
% \operatorname{tr}(H_{1,t}\Sigma_t)
% \big] dt .
% \end{equation}

% Moreover, this quantity equals the change of the distributional risk

% \begin{equation}
% \mathbb E[\Delta_{1\rightarrow2}]=
% \mathcal L_{S_1}(P_T)
% -
% \mathcal L_{S_1}(P_0),
% \end{equation}

% where $P_t$ denotes the distribution of $\theta_t$.

% Hence the trajectory based Itô decomposition and the distributional expansion describe the same quantity from two equivalent perspectives.

% % Corollary (Flatness–Forgetting Relation)
% \begin{Corollary}
%     Under the local quadratic approximation,

% \begin{equation}
% \operatorname{tr}(H_{1,t}\Sigma_t)

% 2,\mathrm{RS}_{S_1}^{(0)}(\theta_t,\Sigma_t).
% \end{equation}

% Therefore the diffusion component of forgetting can be written as

% \begin{equation}
% \frac12
% \int_0^T
% \operatorname{tr}(H_{1,t}\Sigma_t) dt

% \int_0^T
% \mathrm{RS}_{S_1}^{(0)}(\theta_t,\Sigma_t) dt .
% \end{equation}

% \end{Corollary}


% This shows that expected sharpness measures the instantaneous diffusion induced forgetting along the optimization trajectory.

% Consequently, flat solutions characterized by small expected sharpness reduce the accumulation of forgetting during training.




% Let $P=\mathcal{N}(\theta,\Sigma)$
% % \begin{equation}
% % P=\mathcal{N}(\theta,\Sigma)
% % \end{equation}
% denote a Gaussian parameter distribution centered at $\theta$.
% The distributional risk on dataset $S$ is

% \begin{equation}
% \mathcal{L}_S(P)=\mathbb{E}_{W\sim P}\hat L_S(W).
% \end{equation}

% Using a second-order Taylor expansion of $\hat L_S(W)$ around $\theta$ gives

% \begin{equation}
% \mathcal{L}_S(P)
% \approx
% \hat L_S(\theta)
% +
% \frac12 tr(H_S(\theta)\Sigma).
% \end{equation}

% The additional term $tr(H_S(\theta)\Sigma)$
% % \begin{equation}
% % tr(H_S(\theta)\Sigma)
% % \end{equation}
% quantifies the sensitivity of the loss to parameter perturbations.
% This quantity is commonly interpreted as \textbf{expected sharpness} to quantifies the sensitivity of $\hat L_{S}(\theta)$ to a prescribed neighborhood around a point $\theta$.
% For a Gaussian perturbation $\varepsilon\sim\mathcal N(0,\Sigma)$, we define the randomized sharpness
% \begin{equation}
% \mathrm{RS}^{(0)}_{S}(\theta,\Sigma)
% :=
% \mathbb E_{\varepsilon\sim\mathcal N(0,\Sigma)}
% \big[\hat L_{S}(\theta+\varepsilon)\big]
% -
% \hat L_{S}(\theta) = \mathcal{L}_S(P) - \hat L_{S}(\theta)
% \label{eq:random_sharpness_def}
% \end{equation}
% which measure the local character.
% Under the same local quadratic model, the sharpness proxy reduces to the trace form
% \begin{equation}
% \mathrm{RS}^{(0)}_{S}(\theta,\Sigma)
% \approx
% \frac12\,\operatorname{tr}(H_S(\theta)\Sigma).
% \label{eq:sharpness_trace_proxy}
% \end{equation}


% Let $ P_1=\mathcal N(\theta_1,\Sigma_1)$ and $P_2=\mathcal N(\theta_2,\Sigma_2)$ denote the analysis posteriors after tasks 1 and 2.
% We define distributional forgetting on $S_1$ as
% \begin{equation}
% \Delta_{1\to2}^{\mathrm{dist}}
% := \mathcal{L}_S(P_1) - \mathcal{L}_S(P_2)
% = 
% \mathbb E_{w\sim\mathcal N(\theta_1,\Sigma_1)}\hat L_{S_1}(w) - 
% \mathbb E_{w\sim\mathcal N(\theta_2,\Sigma_2)}\hat L_{S_1}(w)
% % \bar L_{S_1}(\theta_2,\Sigma_2)
% % -
% % \bar L_{S_1}(\theta_1,\Sigma_1).
% \label{eq:dist_forgetting_def}
% \end{equation}
% Applying this expansion to the old task yields the distributional forgetting approximation

% \begin{equation}
% \begin{aligned}
%     \Delta_{1 \rightarrow 2}^{\mathrm{dist}}
% &= \underbrace{\hat{L}_{S_1}\left(\theta_2\right)-\hat{L}_{S_1}\left(\theta_1\right)}_{\text{ mean drift}}+\underbrace{\operatorname{RS}_{S_1}^{(0)}\left(\theta_2, \Sigma_2\right)-\operatorname{RS}_{S_1}^{(0)}\left(\theta_1, \Sigma_1\right)}_{\text{sharpness }} 
% % \\
% % & = 
% % \underbrace{g_1^\top\delta_{1\rightarrow2}
% % +
% % \frac12
% % \delta_{1\rightarrow2}^\top H_1\delta_{1\rightarrow2}}_{\text{ drift}}
% % + \underbrace{\frac{1}{2} \operatorname{tr}\left(H_1\left(\Sigma_2-\Sigma_1\right)\right)}_{\text{covariance shift}} +\underbrace{ \frac{1}{2} \operatorname{tr}\left(\left(H_2-H_1\right) \Sigma_2\right)}_{\text{curvature drift}} 
% \end{aligned}
% \end{equation}. The first two terms capture changes in the mean parameter, while the final two term reflects how the covariance of the parameter distribution affects forgetting.


The two analyses above provide complementary perspectives on forgetting.
The Itô-based derivation characterizes the change of the old-task loss as a path integral accumulated along the stochastic optimization trajectory, while the distributional expansion expresses the same quantity in terms of endpoint statistics of the parameter distribution.
More precisely, the drift integral corresponds to the change of the mean loss, whereas the diffusion integral corresponds to the change in sharpness between the two parameter distributions.
Hence the two formulations describe the same underlying quantity from different viewpoints: the stochastic-process analysis provides a \emph{path-integral view} of how forgetting accumulates during training, while the distributional analysis provides an \emph{endpoint view} based on the statistics of the final parameter distribution.
Both reveal that \textbf{forgetting arises from the interaction between optimization dynamics and curvature-sensitive directions of the old task}.

% The two analyses above provide complementary perspectives on forgetting.
% The Itô-based derivation characterizes the change of the old-task loss
% as a path integral along the stochastic optimization trajectory,
% while the distributional expansion expresses the same quantity in terms
% of endpoint statistics of the parameter distribution.
% Hence the two formulations correspond to the same underlying mechanism,
% with the stochastic-process analysis providing a \emph{path-integral view}
% and the distributional analysis providing an \emph{endpoint view}.
% Both reveal that \textbf{forgetting arises from the interaction between
% optimization dynamics and curvature-sensitive directions of the old task}.
% More precisely, the path-integral decomposition
% % \[
% % \Delta L_1
% % =
% % -\int_0^T \nabla L_1(\theta_t)^\top q_t dt
% % +
% % \frac12 \int_0^T tr(H_{1,t}\Sigma_t)dt
% % \]
% describes how drift and diffusion accumulate along the training trajectory.
% By contrast, the distributional expansion
% % \[
% % \Delta_{1\rightarrow2}^{dist}
% % =
% % g_1^\top\delta_{1\rightarrow2}
% % +
% % \frac12\delta_{1\rightarrow2}^\top H_1\delta_{1\rightarrow2}
% % +
% % \frac12 tr(H_1(\Sigma_2-\Sigma_1))
% % +
% % \frac12 tr((H_2-H_1)\Sigma_2)
% % \]
% expresses the same quantity using endpoint statistics of the parameter
% mean and covariance.


\subsection{Cross-Task Noise--Curvature Interaction}

The diffusion term appearing in Theorem~1
\begin{equation}
\operatorname{tr}(H_{1,t}\Sigma_{S_2}^{alg})
\end{equation}
captures the interaction between stochastic perturbations during training and the curvature of the old-task loss.
Importantly, in continual learning these two quantities originate from different tasks.
The covariance $\Sigma_{S_2}^{alg}$ reflects the stochastic perturbations introduced by the optimization process while learning the new task $S_2$, whereas the curvature matrix $H_{1,t}$ characterizes the local sensitivity of the loss associated with the previous task $S_1$.

Consequently, the diffusion component of forgetting arises from a cross-task interaction: stochastic updates generated while optimizing the new task are filtered through the curvature structure of the old-task loss.
To better understand this interaction, consider the eigendecomposition
\begin{equation}
H_{1,t} = U \Lambda U^\top .
\end{equation}

Expressing the noise covariance in the same basis $\tilde{\Sigma}_t = U^\top \Sigma_t U$ yields
\begin{equation}
\operatorname{tr}(H_{1,t}\Sigma_t)
=
\sum_i
\lambda_i \tilde{\sigma}_i .
\end{equation}

This expression shows that stochastic perturbations along directions with large curvature of the old-task loss contribute disproportionately to forgetting.
Even small perturbations can significantly increase the loss of the previous task if they align with highly sensitive directions.

This phenomenon differs fundamentally from flatness analysis in standard optimization.
Existing sharpness-aware methods primarily aim to reduce the curvature of the current-task loss $H_{2}$, thereby promoting flat minima for the task being optimized.
Regularization-based continual learning methods, on the other hand, focus on constraining the deterministic update direction in order to reduce gradient interference.
However, neither class of methods explicitly controls the interaction between the stochastic perturbations generated during new-task training and the curvature structure of the previous tasks.

The term $\operatorname{tr}(H_{1,t}\Sigma_t)$ therefore reveals a distinct mechanism of forgetting: stochastic exploration induced by the new task can amplify the loss of previous tasks when the noise covariance aligns with their sensitive directions.
Understanding and controlling this cross-task noise--curvature interaction is therefore crucial for mitigating forgetting.


\section{Method: Drift–Diffusion Controlled Continual Learning}
\subsection{Constrained Optimization Objective}

Section~3 shows that forgetting during task-$2$ training can be characterized by two mechanisms: gradient interference and curvature-weighted stochastic perturbations. 
In particular, Theorem~1 shows that the expected forgetting from task $1$ to task $2$ contains gradient interference between tasks and  curvature-weighted amplification of stochastic perturbations.

% The first term captures gradient interference between tasks, while the second term measures curvature-weighted amplification of stochastic perturbations.

Our goal is to minimize forgetting while ensuring sufficient learning progress on the new task. 
% From the stochastic optimization dynamics in Eq.~\eqref{eq:noisy_gd_sde_explain}, the learning progress on task $2$ satisfies
% \begin{equation}
% \mathbb{E}\!\left[L_{S_2}(\theta_t)-L_{S_2}(\theta_0)\right]
% =
% -
% \int_0^t
% \mathbb{E}[q_t^\top q_t]dt
% +
% \frac12
% \int_0^t
% \mathbb{E}\!\left[
% \operatorname{tr}(H_{2,t}\Sigma_t)
% \right]dt ,
% \label{eq:task2_progress_method}
% \end{equation}
% where $H_{2,t}=\nabla^2 L_{S_2}(\theta_t)$.
Combining Eqs.~\eqref{eq: current task energy} and 
\eqref{eq: current task forgetting}, we formulate continual learning as a constrained optimization problem over the stochastic training dynamics:

\begin{equation}
\min_{q_t,\Sigma_t}
\;
\int_0^T
\left(
-
v_t^\top q_t
+
\frac12
\operatorname{tr}(H_{1,t}\Sigma_t)
\right)dt
\end{equation}

subject to

\begin{equation}
\int_0^T
\left(
q_t^\top q_t
-
\frac12
\operatorname{tr}(H_{2,t}\Sigma_t)
\right)dt
\ge \kappa ,
\end{equation}

where $\kappa>0$ specifies the minimum learning progress required on task $2$.

This formulation explicitly controls two aspects of the optimization dynamics:

\begin{itemize}
\item \textbf{Drift control:} regulating the update direction $q_t$ to reduce gradient interference with the old-task loss.
\item \textbf{Diffusion control:} shaping the noise covariance $\Sigma_t$ to limit curvature-weighted perturbations on the old-task loss.
\end{itemize}

In the following sections we derive practical update rules that approximately solve this constrained optimization problem during training.






% \section{Method: Drift–Diffusion Controlled Continual Learning}

% Section 3 shows that forgetting arises from two sources in the stochastic optimization dynamics: drift interference and diffusion interference.  
% Based on this analysis, we propose a dual-channel control framework that explicitly regulates both components during training.



% \subsection{Optimization Control Objective}

% Theorem~1 shows that the increase of the old-task loss during new-task training admits the decomposition.
% % \begin{equation}
% %     \mathbb{E}[\Delta_{1\rightarrow 2}]=
% % \mathbb{E}[L_{S_1}(\theta_t)-L_{S_1}(\theta_0)]
% % =
% % - \underbrace{\int_0^T
% % \mathbb{E}[v_t^\top q_t]dt}_{\text{drift interference}}
% % +
% % \underbrace{\frac12
% % \int_0^T
% % \mathbb{E}[tr\left(H_{1,t}\Sigma^{alg}_{S_2}(\theta_t)\right)]dt}_{\text{diffusion interference}} .
% % \end{equation} 
% % \begin{equation}
% % \mathbb{E}[L_{S_1}(\theta_2)-L_{S_1}(\theta_1)]
% % =
% % -\int_0^T
% % \mathbb{E}[\nabla L_{S_1}(\theta_t)^\top q_t]dt
% % +
% % \frac12
% % \int_0^T
% % \mathbb{E}[\operatorname{tr}(H_{1,t}\Sigma_t)]dt .
% % \end{equation}
% The first term measures the interference between the deterministic update direction and the gradient of the old-task loss, while the second term measures the interaction between stochastic perturbations and the curvature of the old-task loss.

% This decomposition suggests a direct control objective for continual learning: minimizing the increase of the old-task loss by regulating both the drift and diffusion components of the optimization dynamics.
% % Assuming stochastic training dynamics of the form
% % \begin{equation}
% % d\theta_t = -q_t dt + \Sigma_t^{1/2} dW_t ,
% % \end{equation}
% we formulate the followig objective
% \begin{equation}
% \min_{q_t,\Sigma_t}
% \int_0^T
% \left(
% \nabla L_{S_1}(\theta_t)^\top q_t
% +
% \frac12
% \operatorname{tr}(H_{1,t}\Sigma_t)
% \right) dt .
% \end{equation}

% The first term corresponds to deterministic gradient interference, while the second term corresponds to curvature-weighted stochastic exploration introduced by the optimizer.
% This objective reveals two complementary control channels:
% (1) regulating the deterministic update direction to reduce gradient interference, and
% (2) shaping the geometry of stochastic perturbations to avoid high-curvature directions of the old-task loss.

% \subsection{Drift Control: Gradient Projection}

\subsection{Drift Control via Gradient Projection}

The analysis in Section~3 shows that the deterministic component of forgetting is governed by the \emph{drift interference} term
\begin{equation}
- v_t^\top q_t ,
\end{equation}
where $v_t=\nabla L_{S_1}(\theta_t)$ denotes the gradient of the old-task loss and $q_t$ is the update direction used for optimizing the new task.

To reduce the increase of the old-task loss, the update direction should avoid moving along directions that are negatively aligned with $v_t$. A natural solution is to remove from the current-task gradient the component that interferes with the old-task gradient~\cite{lopez-pazGradientEpisodicMemory2017,farajtabarOrthogonalGradientDescent2020}.

Let $g_t=\nabla L_{S_2}(\theta_t)$
% \begin{equation}
% g_t=\nabla L_{S_2}(\theta_t)
% \end{equation}
denote the gradient of the current task.
Following the orthogonal gradient descent principle, we construct the update direction as
\begin{equation}
q_t
=
g_t
-
\frac{g_t^\top v_t}{\|v_t\|^2} v_t .
\end{equation}

This projection removes the component of $g_t$ along the old-task gradient and ensures
\begin{equation}
v_t^\top q_t = 0 ,
\end{equation}
thereby eliminating the drift interference term in the forgetting decomposition.

% \paragraph{Maintaining old-task information.}
Computing the projection requires access to gradient information from previous tasks. In practice, this information can be maintained using standard continual learning strategies. One option is to store a small episodic memory of past samples and compute their gradients, as in Gradient Episodic Memory (GEM)~\cite{lopez-pazGradientEpisodicMemory2017}. Another option is to maintain a set of representative historical gradients or a gradient subspace, as in Orthogonal Gradient Descent (OGD)~\cite{farajtabarOrthogonalGradientDescent2020}. These mechanisms provide estimates of the old-task gradient $v_t$ required for the projection step.

% \paragraph{Relation to existing methods.}

% Projection-based updates such as OGD and GEM regulate the update direction to prevent increases in the losses of previous tasks. 
% In those methods the projection constraint is introduced heuristically based on gradient compatibility.
% In contrast, 
Our formulation derives the projection rule directly from the forgetting dynamics established in Section~3. In particular, the projection can be interpreted as an optimal \emph{drift controller} that minimizes the deterministic component of forgetting characterized by the term $-v_t^\top q_t$.
Importantly, our framework is compatible with existing projection-based implementations. The same memory mechanisms used in GEM or the gradient subspace maintained in OGD can be directly integrated into our algorithm to estimate $v_t$.

% \paragraph{Parameter update.}

% The resulting parameter update for task $2$ becomes
% \begin{equation}
% \theta_{t+1}
% =
% \theta_t
% -
% \eta q_t.
% \end{equation}
% In the next subsection we introduce a complementary mechanism that controls the stochastic component of the optimization dynamics.







% The drift component of forgetting arises from the alignment between the deterministic update direction and the gradient of the old-task loss.

% Let

% \begin{equation}
% g_t = \nabla L_{S_2}(\theta_t)
% \end{equation}

% denote the gradient computed from the current minibatch of the new task.

% To identify directions that are sensitive for previous tasks, we estimate a sensitivity subspace using a replay buffer $\mathcal{M}$.  
% Let

% \begin{equation}
% S = \mathrm{span}\{\nabla L(x_i)\}_{x_i\in\mathcal{M}}
% \end{equation}

% denote the subspace spanned by gradients of replay samples.

% Let $B$ be an orthonormal basis of this subspace and define the orthogonal projection

% \begin{equation}
% P_\perp = I - BB^\top .
% \end{equation}

% The controlled update direction is defined as

% \begin{equation}
% q_t = P_\perp g_t .
% \end{equation}

% This projection removes gradient components that lie inside the replay sensitivity subspace, thereby reducing the inner product

% \[
% \nabla L_{S_1}(\theta_t)^\top q_t
% \]

% which corresponds to the drift interference term in Eq.~(13).
\subsection{Diffusion Control: Curvature-Guided Structured Noise}

The analysis in Section~3 shows that stochastic perturbations $\epsilon_t$ influence forgetting through the curvature-weighted quantity $\operatorname{tr}(H_{1,t}\Sigma_t)$.
Therefore, instead of promoting flatness indiscriminately for the current task, we explicitly regulate the diffusion geometry so that stochastic perturbations avoid directions with large curvature in the old-task loss while remaining beneficial for optimizing the new task.

Specifically, stochastic optimization contributes to the progress of the current task through the diffusion term $\frac12 \operatorname{tr}(H_{2,t}\Sigma_t)$,
% \begin{equation}
% \frac12 \operatorname{tr}(H_{2,t}\Sigma_t),
% \end{equation}
while it contributes to forgetting through $\operatorname{tr}(H_{1,t}\Sigma_t)$.
% \begin{equation}
% \operatorname{tr}(H_{1,t}\Sigma_t).
% \end{equation}
This suggests that the noise covariance should satisfy two complementary objectives:
(i) reducing perturbations along directions that are highly sensitive for previous tasks, and
(ii) maintaining sufficient exploration aligned with the geometry of the current task.

% \paragraph{Fisher approximation.}



% \paragraph{Diffusion objective.}
% Under the Fisher approximation, the diffusion control problem becomes
% \begin{equation}
% \min_{\Sigma_t}
% \operatorname{tr}(F_{1,t}\Sigma_t)
% \end{equation}
% subject to maintaining exploration useful for learning the new task
% \begin{equation}
% \operatorname{tr}(F_{2,t}\Sigma_t) \ge \gamma ,
% \end{equation}
% together with a fixed noise energy constraint
% \begin{equation}
% \operatorname{tr}(\Sigma_t)=\sigma^2 .
% \end{equation}

This leads to the constrained optimization problem
\begin{equation}
\min_{\Sigma_t}
\operatorname{tr}(H_{1,t}\Sigma_t)
\quad
\text{s.t.}
\quad
\operatorname{tr}(H_{2,t}\Sigma_t)\ge\gamma,
\quad
\operatorname{tr}(\Sigma_t)=\sigma^2 .
\label{eq: diffusion target}
\end{equation}
This formulation explicitly balances two goals: minimizing the curvature-weighted perturbation on the old-task loss while preserving stochastic exploration aligned with the geometry of the current task.
% \paragraph{Closed-form diffusion controller.}
The constrained problem above admits an analytical solution.
% Let $F_{1,t}$ and $F_{2,t}$ denote the Fisher information matrices of the old and current tasks respectively.




\begin{proposition}
Under the Lagrangian relaxation of the diffusion objective in Eq.~\eqref{eq: diffusion target}, i.e.,\ $\min_{\Sigma_t} \operatorname{tr}(H_{1,t}\Sigma_t) - \beta\,\operatorname{tr}(H_{2,t}\Sigma_t)$ subject to $\operatorname{tr}(\Sigma_t)=\sigma^2$ and $\Sigma_t\succeq 0$, the optimal covariance matrix $\Sigma_t^\star$ is

\begin{equation}
\Sigma_t^\star
=
\sigma^2
\bigl(F_{1,t} - \beta F_{2,t} + \varepsilon I\bigr)^{-1},
\label{eq:diffusion_controller}
\end{equation}

where $F_{1,t}$ and $F_{2,t}$ are the Fisher information matrices for the old and current tasks, $\sigma^2$ controls the total noise energy, $\beta\geq 0$ balances the trade-off between avoiding old-task curvature and encouraging new-task exploration, and $\varepsilon>0$ ensures positive definiteness. The proof is provided in Appendix~A.
\end{proposition}

Equation~(\ref{eq:diffusion_controller}) reveals that the optimal noise covariance suppresses perturbations along directions where $F_{1,t}$ is large (high curvature for old tasks) and amplifies them along directions where $F_{2,t}$ is large (informative for the current task). In practice we use the conservative special case $\beta=0$, setting $\Sigma_t^\star = \sigma^2(F_{1,t}+\varepsilon I)^{-1}$, which avoids old-task curvature without requiring an online estimate of $F_{2,t}$.
This structured noise therefore acts as a \emph{diffusion controller} that complements the drift projection mechanism in Section~4.2. 




Directly computing the Hessian matrices is typically intractable for modern neural networks.
Following common practice in continual learning and Bayesian approximations, we approximate the curvature using the Fisher information matrices
\begin{equation}
H_{1,t} \approx F_{1,t}, 
\qquad
H_{2,t} \approx F_{2,t}.
\end{equation}
The Fisher matrix captures the sensitivity of the task likelihood with respect to parameter perturbations and can be estimated using gradients from stored samples or historical statistics.
% 22222222222222
%%%%%%%%

% \subsection{Diffusion Control: Curvature-Guided Structured Noise}

% The analysis in Section~3 shows that stochastic perturbations $\epsilon_t$ influence forgetting through the curvature-weighted quantity $\operatorname{tr}(H_{1,t}\Sigma_t)$.
% % \begin{equation}
% % \operatorname{tr}(H_{1,t}\Sigma_t).
% % \end{equation}
% Therefore, instead of promoting flatness of current task $\operatorname{tr}(H_{2,t}\Sigma_t)$, we directly regulate the diffusion geometry so that stochastic perturbations avoid directions with large curvature in the old-task loss.

% Specifically, we construct a noise covariance matrix that reduces perturbations along high-curvature directions of $H_{1,t}$ while preserving exploration in other directions. This mechanism acts as a \emph{diffusion controller} that complements the drift projection in Section~4.2.

% % \paragraph{Closed-form diffusion controller.}

% We now derive the optimal diffusion structure under the objective in Section~4.1.
% Recall that stochastic perturbations influence forgetting through the curvature-weighted quantity

% \begin{equation}
% \operatorname{tr}(H_{1,t}\Sigma_t).
% \end{equation}

% Minimizing this term directly requires controlling the covariance matrix of the optimization noise.
% However, computing the full Hessian of the old-task loss is typically intractable for large neural networks.
% Following common practice in continual learning and Bayesian approximations, we adopt the Fisher information matrix as a curvature proxy $H_{1,t} \approx F_{1,t}$.
% % \begin{equation}
% % H_{1,t} \approx F_{1,t}.
% % \end{equation}
% The Fisher matrix captures the sensitivity of the old-task likelihood with respect to parameter perturbations and can be estimated using reserved gradients or historical statistics. 

% Substituting this approximation into the diffusion objective yields
% \begin{equation}
% \min_{\Sigma_t}
% \operatorname{tr}(F_{1,t}\Sigma_t).
% \end{equation}
% To avoid collapsing the stochastic exploration entirely, we impose a fixed noise energy constraint
% \begin{equation}
% \operatorname{tr}(\Sigma_t) = \sigma^2 .
% \end{equation}
% This leads to the constrained optimization problem
% \begin{equation}
% \min_{\Sigma_t}
% \operatorname{tr}(F_{1,t}\Sigma_t)
% \quad
% \text{s.t.}
% \quad
% \operatorname{tr}(\Sigma_t)=\sigma^2 .
% \end{equation}
% 333333333333333
% The solution is obtained by allocating smaller variance along directions with large Fisher curvature.
% Let

% \begin{equation}
% F_{1,t}=U\Lambda U^\top
% \end{equation}

% be the eigendecomposition of the Fisher matrix.
% The optimal covariance takes the form

% \begin{equation}
% \Sigma_t
% =
% U
% \operatorname{diag}(\sigma_i^2)
% U^\top
% \end{equation}

% with

% \begin{equation}
% \sigma_i^2 \propto \frac{1}{\lambda_i+\epsilon}.
% \end{equation}

% Equivalently, the covariance can be written in closed form as

% \begin{equation}
% \Sigma_t
% =
% \sigma^2 (F_{1,t}+\epsilon I)^{-1}.
% \end{equation}

% This structure reduces stochastic perturbations along directions that are highly sensitive for previous tasks while preserving exploration in directions that are less critical.

% \paragraph{Final stochastic update.}

% Combining the drift controller in Section~4.2 with the diffusion controller derived above yields the final stochastic update

% \begin{equation}
% \theta_{t+1}
% =
% \theta_t
% -
% \eta q_t
% +
% \eta \Sigma_t^{1/2}\epsilon_t,
% \qquad
% \epsilon_t\sim\mathcal N(0,I).
% \end{equation}

% This update simultaneously regulates

% \begin{itemize}
% \item deterministic gradient interference through the projection step, and
% \item stochastic perturbation geometry through curvature-aware diffusion.
% \end{itemize}

% Consequently, the optimization dynamics directly minimize both components of the forgetting decomposition derived in Section~3.







% 333333333333333333333333333


% The diffusion component of forgetting arises from the interaction between optimization noise and the curvature of the old-task loss, quantified by the term

% \[
% \operatorname{tr}(H_{1,t}\Sigma_t).
% \]

% Section 3.3 shows that forgetting is amplified when stochastic perturbations align with directions of high curvature of the old-task loss.

% To reduce this interaction, we design the noise covariance so that perturbations occur preferentially in directions where the old-task loss is relatively flat.

% For a fixed noise magnitude constraint, minimizing the diffusion term leads to an optimal covariance structure proportional to

% \begin{equation}
% \Sigma^\star \propto H_{1,t}^{-1}.
% \end{equation}

% Computing the full Hessian is impractical in deep networks.  
% Following common practice, we approximate the curvature using the empirical Fisher matrix

% \begin{equation}
% H_{1,t} \approx F_{1,t}.
% \end{equation}

% To maintain consistency with the drift controller, we restrict the covariance to the orthogonal subspace defined by $P_\perp$.  
% The resulting covariance approximation becomes

% \begin{equation}
% \Sigma^\star \approx P_\perp (F_{1,t}+\epsilon)^{-1} P_\perp .
% \end{equation}

% To generate perturbations with this covariance, we sample

% \[
% z \sim \mathcal N(0,I)
% \]

% and transform it as

% \begin{equation}
% \xi_t = P_\perp (F_{1,t}+\epsilon)^{-1/2} P_\perp z .
% \end{equation}

% This transformation injects larger variance along directions where the old-task loss is flat while suppressing perturbations along sensitive directions.
\subsection{Unified Drift–Diffusion Update}

Sections~4.2 and 4.3 derive two complementary controllers for continual learning:

\begin{itemize}
\item \textbf{Drift controller}: projection-based updates that regulate the descent direction $q_t$ to reduce gradient interference with previous tasks.
\item \textbf{Diffusion controller}: structured stochastic perturbations with covariance $\Sigma_t$ that suppress curvature-weighted perturbations on the old-task loss while preserving exploration aligned with the new task.
\end{itemize}

Combining the two components yields the following stochastic update rule

\begin{equation}
\theta_{t+1}
=
\theta_t
-
\eta q_t
+
{\eta}\,\epsilon_t , \quad \epsilon_t \sim \mathcal N(0,\Sigma^*_t)
\end{equation}

where $q_t
=
g_t
-
\frac{g_t^\top v_t}{\|v_t\|^2} v_t$ denotes the drift-controlled gradient obtained via projection (Section~4.2), and $\Sigma_t^\star
=
\sigma^2
\left(
M_t + \epsilon I
\right)^{-1}$
% \begin{equation}
% \epsilon_t \sim \mathcal N(0,\Sigma_t)
% \end{equation}
is the structured perturbation derived in Section~4.3.


\paragraph{Orthogonal noise injection.}

A potential concern is that adding stochastic perturbations may alter the descent direction of the gradient update.
To ensure that the diffusion term only affects the exploration geometry without biasing the deterministic descent direction, we restrict the injected noise to the subspace orthogonal to the current gradient.

Let $q_t$ denote the modified gradient of the current task.
Given a sampled perturbation $\epsilon_t \sim \mathcal N(0,\Sigma_t^*)$, we construct its orthogonal component as
\begin{equation}
\epsilon_t^\perp
=
\epsilon_t
-
\frac{q_t^\top \epsilon_t}{\|q_t\|^2}\, q_t .
\end{equation}
The resulting parameter update becomes
\begin{equation}
\theta_{t+1}
=
\theta_t
-
\eta q_t
+
{\eta}\,\epsilon_t^\perp .
\end{equation}

By construction, the injected perturbation satisfies $g_t^\top \epsilon_t^\perp = 0 ,$
% \begin{equation}
% g_t^\top \epsilon_t^\perp = 0 ,
% \end{equation}
which guarantees that the stochastic perturbation does not modify the descent direction.
Instead, it only broadens the optimization trajectory within the subspace orthogonal to the gradient.

% In practice, computing the exact orthogonal projection can be simplified by approximating the orthogonal component using the cosine similarity between the gradient and the sampled perturbation:

% \begin{equation}
% \epsilon_t^\perp
% \approx
% \epsilon_t
% -
% \sigma g_t ,
% \end{equation}

% where $\sigma = \cos(g_t,\epsilon_t)$ denotes their cosine similarity.
% Following prior work, $\sigma$ can be treated as a constant during optimization for computational efficiency.

This orthogonalization step ensures that the proposed diffusion controller modifies the geometry of stochastic exploration while preserving the deterministic descent dynamics.

The resulting update simultaneously regulates the deterministic descent direction and the stochastic exploration geometry, thereby controlling both sources of forgetting identified in Section~3.

% 6666666

% \subsection{Unified Drift–Diffusion Update}

% Together, the two components form a unified optimization strategy:

% \begin{itemize}
% \item \textbf{Drift control}: projection-based updates reduce gradient interference.
% \item \textbf{Diffusion control}: curvature-aware perturbations reduce noise-induced forgetting.
% \end{itemize}

% The resulting training dynamics simultaneously regulate both deterministic and stochastic components of optimization, directly minimizing the forgetting dynamics derived in Section~3.
% Combining the drift and diffusion controllers yields the following update rule

% \begin{equation}
% \theta_{t+1}
% =
% \theta_t
% -
% \eta P_\perp g_t
% +
% \sqrt{\eta}\sigma
% P_\perp(F_{1,t}+\epsilon)^{-1/2}P_\perp z .
% \end{equation}

% The first term regulates deterministic updates to avoid directions that increase the old-task loss, while the second term shapes stochastic exploration so that perturbations occur primarily in directions where the old-task loss is insensitive.

% Algorithm~1 summarizes the resulting training procedure.

% Overall, the proposed method directly implements the control objective derived in Section~3 by simultaneously regulating the deterministic drift and stochastic diffusion of the optimization dynamics.

% %---------------------------------
% % \section{Method: Drift–Diffusion Dual-Channel Control}

% % \subsection{Control Objective}

% % The analyses above reveal that forgetting during task-2 training is governed by two mechanisms.
% % Theorem~1 shows that the expected increase of the old-task loss admits the decomposition

% % \begin{equation}
% % \mathbb{E}[L_{S_1}(\theta_2)-L_{S_1}(\theta_1)]
% % =
% % -\int_0^T
% % \mathbb{E}[\nabla L_{S_1}(\theta_t)^\top q_t]dt
% % +
% % \frac12
% % \int_0^T
% % \mathbb{E}[\operatorname{tr}(H_{1,t}\Sigma_t)]dt .
% % \end{equation}

% % The first term captures the deterministic drift of the optimization trajectory and measures the interference between the update direction and the gradient of the old-task loss.
% % The second term captures the diffusion effect and quantifies how stochastic perturbations interact with the curvature of the old-task loss.

% % These observations suggest a direct optimization objective for continual learning: controlling both the drift and diffusion components of the training dynamics in order to minimize the increase of the old-task loss.

% % Consider stochastic training dynamics of the form

% % \begin{equation}
% % d\theta_t
% % =
% % -q_t dt
% % +
% % \Sigma_t^{1/2} dW_t .
% % \end{equation}

% % Our objective is to minimize the expected forgetting accumulated along the optimization trajectory

% % \begin{equation}
% % \min_{q_t,\Sigma_t}
% % \int_0^T
% % \left(
% % \nabla L_{S_1}(\theta_t)^\top q_t
% % +
% % \frac12
% % \operatorname{tr}(H_{1,t}\Sigma_t)
% % \right)dt .
% % \end{equation}

% % This formulation reveals two complementary control channels:

% % \begin{itemize}
% % \item \textbf{Drift control}, which regulates the update direction $q_t$ to reduce interference with the gradient of the old task.

% % \item \textbf{Diffusion control}, which shapes the noise covariance $\Sigma_t$ to reduce stochastic exploration along high-curvature directions of the old-task loss.
% % \end{itemize}

% % This perspective provides a unified view of existing continual learning strategies.
% % Regularization and gradient-projection methods primarily influence the drift component, whereas flatness-aware optimization methods modify the diffusion component by reshaping the effective noise geometry during training.

% % In the next section we introduce an algorithm that jointly controls both mechanisms.

% % Section~3 shows that the increase of the old-task loss during new-task training arises from two components of the optimization dynamics: drift interference and diffusion interference. 
% % The first component is determined by the alignment between the deterministic update direction and the gradient of the old-task loss, while the second component is determined by the interaction between the optimizer noise covariance and the curvature of the old-task loss.

% % Motivated by this analysis, we propose a dual-channel control framework that regulates both components of the optimization dynamics during continual learning.
% % The framework consists of two complementary mechanisms:

% % \begin{itemize}
% % \item a \textbf{drift controller} that suppresses update directions which directly increase the old-task loss,
% % \item a \textbf{diffusion controller} that shapes the geometry of stochastic perturbations to avoid high-curvature directions of the old-task loss.
% % \end{itemize}

% % Together these two components implement the control objective derived in Section~3.4 by simultaneously regulating deterministic updates and stochastic exploration during training.



% % \subsection{Drift Channel: Orthorgng Gradient Projection}

% % The drift channel controls how deterministic updates influence previous tasks.

% % Let
% % \begin{equation}
% % g_t = \nabla L_t(\theta_t)
% % \end{equation}

% % denote the gradient computed from the current minibatch of the new task.

% % To estimate directions that are sensitive for previous tasks, we compute gradients using a replay buffer $\mathcal{M}$. 
% % The sensitivity subspace is defined as

% % \begin{equation}
% % \mathcal{S}
% % =
% % \operatorname{span}
% % \{\nabla L(x_i)\}_{x_i \in \mathcal{M}} .
% % \end{equation}

% % Let $B$ denote an orthonormal basis of this subspace. 
% % The orthogonal projection matrix is

% % \begin{equation}
% % P_\perp
% % =
% % I - BB^\top .
% % \end{equation}

% % The drift controller removes gradient components lying inside this sensitive subspace.
% % The controlled update direction is therefore

% % \begin{equation}
% % q_t
% % =
% % P_\perp g_t .
% % \end{equation}

% % This projection suppresses deterministic updates that would directly increase the loss of previous tasks.



% % \subsection{Diffusion Channel: Curvature-Guided Structured Noise}

% % The diffusion channel controls how stochastic perturbations interact with the curvature of the old-task loss.

% % Section~3 shows that minimizing the curvature–noise interaction term
% % $\operatorname{tr}(H_1 \Sigma)$ leads to an optimal covariance structure of the form

% % \begin{equation}
% % \Sigma^\star
% % =
% % P_\perp H_1^{-1} P_\perp ,
% % \end{equation}

% % where $H_1$ denotes the Hessian of the old-task loss.

% % Computing the full Hessian is infeasible in deep networks.
% % Following common practice, we approximate the curvature using the empirical Fisher matrix

% % \begin{equation}
% % H_1 \approx F_1 .
% % \end{equation}

% % Using a diagonal approximation yields

% % \begin{equation}
% % \Sigma^\star
% % \approx
% % P_\perp
% % \operatorname{diag}(F_1+\epsilon)^{-1}
% % P_\perp .
% % \end{equation}

% % To generate perturbations with this covariance, we sample

% % \begin{equation}
% % z \sim \mathcal{N}(0,I)
% % \end{equation}

% % and transform it as

% % \begin{equation}
% % \xi_t
% % =
% % P_\perp
% % (F_1+\epsilon)^{-1/2}
% % P_\perp z .
% % \end{equation}

% % This transformation injects larger variance along directions where the old-task loss is relatively flat while suppressing perturbations along sensitive directions.



% % \subsection{Unified Update Rule}

% % Combining the drift and diffusion controllers yields the following update rule

% % \begin{equation}
% % \theta_{t+1}
% % =
% % \theta_t
% % -
% % \eta P_\perp g_t
% % +
% % \sqrt{\eta}\sigma
% % P_\perp
% % (F_1+\epsilon)^{-1/2}
% % P_\perp z .
% % \end{equation}

% % The first term regulates deterministic drift to avoid directions that interfere with previous tasks.
% % The second term shapes stochastic exploration so that perturbations primarily occur in directions where the old-task loss is insensitive.

% % Algorithm~\ref{alg:dual_channel_control} summarizes the resulting training procedure.



% % \begin{algorithm}[t]
% % \caption{Replay-Aware Drift–Diffusion Dual-Channel Control}
% % \label{alg:dual_channel_control}
% % \begin{algorithmic}[1]

% % \STATE \textbf{Input:} task datasets $\{\mathcal S_t\}_{t=1}^{T}$, replay buffer $\mathcal M$
% % \STATE Initialize parameters $\theta$

% % \FOR{each task $t=1,\dots,T$}

% % \FOR{each minibatch $B_{\text{new}}$}

% % \STATE Compute new-task gradient
% % \[
% % g \leftarrow \nabla L(B_{\text{new}})
% % \]

% % \STATE Estimate replay gradient subspace
% % \[
% % B \leftarrow \operatorname{basis}\{\nabla L(x_i)\}_{x_i\in\mathcal M}
% % \]

% % \STATE Construct projection matrix
% % \[
% % P_\perp \leftarrow I - BB^\top
% % \]

% % \STATE Drift control
% % \[
% % q \leftarrow P_\perp g
% % \]

% % \STATE Estimate Fisher matrix $F_{\text{old}}$

% % \STATE Sample Gaussian noise
% % \[
% % z \sim \mathcal N(0,I)
% % \]

% % \STATE Structured noise transformation
% % \[
% % z \leftarrow P_\perp (F_{\text{old}}+\epsilon)^{-1/2} P_\perp z
% % \]

% % \STATE Parameter update
% % \[
% % \theta
% % \leftarrow
% % \theta
% % -
% % \eta q
% % +
% % \sqrt{\eta}\sigma z
% % \]

% % \ENDFOR
% % \ENDFOR

% % \end{algorithmic}
% % \end{algorithm}

\begin{algorithm}[t]
\caption{Drift–Diffusion Controlled Continual Learning (Perturbation Form)}
\label{alg:ddcl}
\begin{algorithmic}[1]

\STATE \textbf{Input:}
model parameters $\theta$,
learning rate $\eta$,
memory buffer $M$ for previous tasks,
noise scale $\sigma$,
balance coefficient $\beta$

\FOR{each training iteration}

\STATE Sample minibatch $B$ from current task

\STATE Compute gradient
\[
g_t = \nabla_\theta L_B(\theta_t)
\]

\STATE \textbf{Drift control}

Project gradient to avoid interference with old tasks
\[
q_t = g_t -
\frac{g_t^\top v_t}{\|v_t\|^2} v_t
\]

where $v_t$ denotes the stored gradient direction from previous tasks.

\STATE \textbf{Diffusion covariance}
Estimate Fisher matrices
\[
F_{1,t} \quad \text{(old tasks)}
\]

% \[
% F_{2,t} \approx g_t g_t^\top
% \]

Construct structured covariance
\[
\Sigma_t =
\sigma^2 (F_{1,t}+\epsilon I)^{-1}
\]

\STATE \textbf{Sample perturbation}
\[
\delta_t \sim \mathcal N(0,\Sigma_t)
\]

\STATE \textbf{Orthogonalize perturbation}
\[
\delta_t^\perp
=
\delta_t -
\frac{q_t^\top \delta_t}{\|q_t\|^2} q_t
\]

This ensures the perturbation does not change the descent direction.

\STATE \textbf{Perturbed gradient}
\[
q'_t = \nabla_\theta L_B(\theta_t+\delta_t^\perp)
\]

\STATE \textbf{Diffusion correction}
\[
g_{corr}
=
q_t + \lambda (q'_t-q_t)
\]
where $\lambda$ controls the strength of diffusion regularization.

\STATE \textbf{Parameter update}
\[
\theta_{t+1}
=
\theta_t
-
\eta g_{corr}
\]

\STATE Store gradient or sample in memory buffer $M$

\ENDFOR

\end{algorithmic}
\end{algorithm}



% \begin{algorithm}[t]
% \caption{Drift–Diffusion Controlled Continual Learning}
% \begin{algorithmic}[1]

% \STATE Input: parameters $\theta$, learning rate $\eta$, memory $M$, noise scale $\sigma$

% \WHILE{not converged}

% \STATE Sample minibatch $B$

% \STATE Compute gradient
% \[
% g_t=\nabla L_B(\theta_t)
% \]

% \STATE \textbf{Drift statistics}

% Retrieve old-task gradient direction
% \[
% v_t
% \]

% \STATE \textbf{Drift control}

% Project gradient
% \[
% q_t =
% g_t -
% \frac{g_t^\top v_t}{\|v_t\|^2}v_t
% \]

% \STATE \textbf{Diffusion statistics}

% Estimate Fisher matrices

% \[
% F_{1,t}, \quad F_{2,t}\approx g_t g_t^\top
% \]

% \STATE \textbf{Perturbation construction}

% \[
% \Sigma_t
% =
% \sigma^2 (F_{1,t}-\beta F_{2,t}+\epsilon I)^{-1}
% \]

% Sample perturbation

% \[
% \delta_t \sim N(0,\Sigma_t)
% \]

% \STATE \textbf{Orthogonalize perturbation}

% \[
% \delta_t^\perp
% =
% \delta_t -
% \frac{q_t^\top \delta_t}{\|q_t\|^2}q_t
% \]

% \STATE \textbf{Perturbed gradient}

% \[
% q'_t = \nabla L_B(\theta_t+\delta_t^\perp)
% \]

% \STATE \textbf{Update}

% \[
% \theta_{t+1}
% =
% \theta_t
% -
% \eta
% (q_t+\lambda(q'_t-q_t))
% \]

% \ENDWHILE

% \end{algorithmic}
% \end{algorithm}



% \begin{algorithm}[t]
% \caption{Drift–Diffusion Controlled Continual Learning}
% \begin{algorithmic}[1]
% \STATE Input: model parameters $\theta$, learning rate $\eta$, Stored task statistics for previous tasks,Noise scale $\sigma$

% \FOR{each training epoch}

% \STATE Sample minibatch $B$ from current task

% \STATE Compute batch’s loss gradient $g_{t} = \nabla_\theta L_{B}(\theta_t)$
% % \[
% % g_{2,t} = \nabla_\theta L_{S_2}(\theta_t)
% % \]

% \STATE \textbf{Drift control}
% \STATE \hspace{0.5cm}
% Project gradient to avoid interference with old tasks
% $q_t
% =
% g_t
% -
% \frac{g_t^\top v_t}{\|v_t\|^2} v_t$
% % \[
% % \tilde g_t = \Pi_{\mathcal S}(g_{2,t})
% % \]

% \STATE \textbf{Diffusion control 1}

% % Compute coverance
% % \[
% % \Sigma_t^*
% % =
% % \sigma^2 (F_{1,t}-\beta F_{2,t}+\epsilon I)^{-1}
% % \]
% \STATE \hspace{0.5cm}
% Sample noise from  $\epsilon_t \sim \mathcal N(0,\Sigma_t^*),\quad \Sigma_t^*
% =
% \sigma^2 (F_{1,t}-\beta F_{2,t}+\epsilon I)^{-1}$

% % \[
% % \epsilon_t \sim \mathcal N(0,\Sigma_t^*)
% % \]
% \STATE  \hspace{0.5cm} noise project  $
% \epsilon_t^\perp
% =
% \epsilon_t
% -
% \frac{q_t^\top \epsilon_t}{\|q_t\|^2}\, q_t .
% $

% % \STATE \textbf{Diffusion control 2}

% % % Compute coverance
% % % \[
% % % \Sigma_t^*
% % % =
% % % \sigma^2 (F_{1,t}-\beta F_{2,t}+\epsilon I)^{-1}
% % % \]
% % \STATE \hspace{0.5cm}
% % $q_t' = \nabla_{\theta}L_{B}(\theta_t + q_t),\quad \epsilon_t = 0$



% % Sample noise from  $\epsilon_t \sim \mathcal N(0,\Sigma_t^*),\quad \Sigma_t^*
% % =
% % \sigma^2 (F_{1,t}-\beta F_{2,t}+\epsilon I)^{-1}$

% % \[
% % \epsilon_t \sim \mathcal N(0,\Sigma_t^*)
% % \]
% % \STATE  \hspace{0.5cm} noise project  $
% % \epsilon_t^\perp
% % =

% % $


% \STATE \textbf{Unified update}
% $
% \theta_{t+1}
% =
% \theta_t
% -
% \eta\,q_t
% +
% {\eta}\,\epsilon_t
% $
% \STATE \textbf{Accumulate}  gradient or sample to $M$ for next task 

% \STATE \hspace{0.5cm}

% \ENDFOR

% \end{algorithmic}
% \end{algorithm}






% \begin{algorithm}[t]
% \caption{Drift–Diffusion Controlled Continual Learning}
% \begin{algorithmic}[1]
% \STATE Input: model parameters $\theta$, learning rate $\eta$, Stored task statistics for previous tasks,Noise scale $\sigma$

% \FOR{each training epoch}

% \STATE Sample minibatch $B$ from current task

% \STATE Compute batch’s loss gradient $g_{t} = \nabla_\theta L_{B}(\theta_t)$
% % \[
% % g_{2,t} = \nabla_\theta L_{S_2}(\theta_t)
% % \]

% \STATE \textbf{Drift control}
% \STATE \hspace{0.5cm}
% Project gradient to avoid interference with old tasks
% $q_t
% =
% g_t
% -
% \frac{g_t^\top v_t}{\|v_t\|^2} v_t$
% % \[
% % \tilde g_t = \Pi_{\mathcal S}(g_{2,t})
% % \]

% % \STATE \textbf{Diffusion control 1}

% % % Compute coverance
% % % \[
% % % \Sigma_t^*
% % % =
% % % \sigma^2 (F_{1,t}-\beta F_{2,t}+\epsilon I)^{-1}
% % % \]
% % \STATE \hspace{0.5cm}
% % Sample noise from  $\epsilon_t \sim \mathcal N(0,\Sigma_t^*),\quad \Sigma_t^*
% % =
% % \sigma^2 (F_{1,t}-\beta F_{2,t}+\epsilon I)^{-1}$

% % % \[
% % % \epsilon_t \sim \mathcal N(0,\Sigma_t^*)
% % % \]
% % \STATE  \hspace{0.5cm} noise project  $
% % \epsilon_t^\perp
% % =
% % \epsilon_t
% % -
% % \frac{q_t^\top \epsilon_t}{\|q_t\|^2}\, q_t .
% % $

% \STATE \textbf{Diffusion control 2}

% % Compute coverance
% % \[
% % \Sigma_t^*
% % =
% % \sigma^2 (F_{1,t}-\beta F_{2,t}+\epsilon I)^{-1}
% % \]
% \STATE \hspace{0.5cm}
% $
% \delta = \rho \frac{q_t}{\|q_t\|_2} \quad
% q_t' = \nabla_{\theta}L_{B}(\theta_t + \delta),\quad (\epsilon_t = q_t-q't)$



% % Sample noise from  $\epsilon_t \sim \mathcal N(0,\Sigma_t^*),\quad \Sigma_t^*
% % =
% % \sigma^2 (F_{1,t}-\beta F_{2,t}+\epsilon I)^{-1}$

% % \[
% % \epsilon_t \sim \mathcal N(0,\Sigma_t^*)
% % \]
% % \STATE  \hspace{0.5cm} noise project  $
% % \epsilon_t^\perp
% % =

% % $


% \STATE \textbf{Unified update}
% $
% \theta_{t+1}=
%  \theta_t -\eta  q'_t
% $
% \STATE \textbf{Accumulate}  gradient or sample to $M$ for next task 

% \STATE \hspace{0.5cm}

% \ENDFOR

% \end{algorithmic}
% \end{algorithm}





\section{Experiments}
\label{sec:experiments}

We evaluate on \textbf{ImageNet-R}~\cite{hendrycks2021many} under the class-incremental protocol with 10 tasks of 20 classes each (200 classes total), following~\cite{bian2024continual}.
The backbone is ViT-B/16 pretrained on ImageNet-21K.
All methods use SGD with learning rate $10^{-3}$, cosine scheduler, 20 epochs per task, and batch size 128.
We report three metrics: FAA (Final Average Accuracy, accuracy over all seen classes after the last task), AAA (Average of Average Accuracy, mean per-task cumulative accuracy as tasks arrive), and BWT (Backward Transfer, average forgetting).

\paragraph{Main results.}
Table~\ref{tab:main} compares our method against OGD variants under different noise and perturbation strategies.
The results directly validate the theoretical predictions of Theorem~1:

\begin{table}[h]
\centering
\caption{ImageNet-R class-incremental, 10 tasks $\times$ 20 classes, seed 1993.
$\uparrow$ higher is better; $\downarrow$ lower is better.}
\label{tab:main}
\begin{tabular}{lcccccc}
\toprule
Method & Drift & Diffusion & FAA$\uparrow$ & AAA$\uparrow$ & BWT$\uparrow$ & NME-AAA$\uparrow$ \\
\midrule
OGD (baseline)~\cite{farajtabarOrthogonalGradientDescent2020} & \checkmark & \texttimes & 23.58 & 40.74 & $-$61.30 & 71.52 \\
OGD + Gaussian Noise & \checkmark & isotropic & 17.0\textsuperscript{$\dagger$} & 33.9\textsuperscript{$\dagger$} & $\downarrow\downarrow$ & 54.4\textsuperscript{$\dagger$} \\
OGD + RWP Gaussian & \checkmark & approx.\ & 36.67 & 46.53 & $-$44.87 & 71.50 \\
OGD + RWP Fisher (ours) & \checkmark & \checkmark Fisher & 35.41 & 48.98 & $-$47.92 & 72.97 \\
SAM-OGD & SAM+proj & \texttimes & 53.12 & 59.25 & $-$27.87 & 74.98 \\
\textbf{GAM-OGD-Fisher (ours)} & \textbf{GAM+proj} & \textbf{\checkmark Fisher} & \textbf{63.61} & \textbf{70.33} & $\mathbf{-22.07}$ & \textbf{79.45} \\
\bottomrule
\end{tabular}
\footnotesize{$\dagger$ CNN per-task log (FAA/AAA not yet in summary sheet).}
\end{table}

\paragraph{Ablation: diffusion interference (Theorem~1 validation).}
\textit{Isotropic noise hurts.}
OGD + Gaussian noise reduces CNN accuracy by $\approx 16.7$ points versus OGD alone, confirming that tr$(H_1\Sigma_\text{iso})$ with $\Sigma_\text{iso}\propto I$ amplifies forgetting when noise aligns with high-curvature directions of old tasks.
\textit{Fisher-guided noise helps.}
OGD + RWP Fisher gains $+8.2$ AAA points over OGD (48.98 vs.\ 40.74), with BWT improving from $-61.3$ to $-47.9$.
\textit{Structure matters.}
RWP Fisher (48.98 AAA) outperforms RWP Gaussian (46.53 AAA), confirming that the Fisher weighting in Eq.~(\ref{eq:diffusion_controller}) is the operative mechanism.
\textit{Combined drift + diffusion control is best.}
GAM-OGD-Fisher achieves AAA$=70.33$ ($+29.6$ vs.\ OGD, BWT from $-61.3$ to $-22.1$), directly validating the joint control objective in Eq.~(\ref{eq:cl_constrained_obj}).


\section{Conclusion}
\label{sec:conclusion}

We have presented a unified framework for continual learning that explicitly controls both the drift and diffusion components of the stochastic optimization dynamics.
Using It\^{o}'s lemma, we derived a decomposition of the expected forgetting into gradient interference and a cross-task noise--curvature interaction term $\operatorname{tr}(H_1\Sigma)$, and showed that existing regularization-based and gradient projection methods address only the drift component, leaving diffusion interference uncontrolled.
Our proposed diffusion controller introduces a structured noise covariance $\Sigma^\star = \sigma^2(F_1 - \beta F_2 + \varepsilon I)^{-1}$ that selectively suppresses perturbations in directions sensitive to old tasks, with $\beta=0$ as a conservative default.
Experiments on ImageNet-R validate all three theoretical predictions: uncontrolled isotropic diffusion amplifies forgetting, Fisher-guided diffusion control reduces it, and the combined drift--diffusion controller achieves strong class-incremental performance.
Future work includes extending the forgetting analysis beyond the two-task case, deriving data-dependent $\beta$ selection from the $\operatorname{tr}(H_2\Sigma)\geq\gamma$ constraint, and scaling the Fisher approximation to LoRA-parameterized models.


\section*{References}

\bibliographystyle{unsrt}
% \bibliographystyle{plainnat}
\bibliography{ref}






% % \section{Bridging Noise Injection and LoRA}
% % \section{Towards Structured Noise Injection for Robust and Generalizable LoRA Adaptation}


% \section{Appendix: Theoretical Proofs}

% % \subsection{Second-order objects: Hessian, GGN, Fisher, and expected Fisher}
% % \label{appendix:second-order-metrics}
% % This appendix summarizes the relations between commonly used second-order quantities and clarifies which object is meant when we write $C_t$.
% % We follow the unified viewpoint advocated in ``New Insights and Perspectives on the Natural Gradient Method''~\citep{JMLR:NewInsights}.

% % \paragraph{Setup.}
% % Consider a (single-task) loss of the form
% % \begin{equation}
% % \ell(W;z)=\ell\big(f_W(x),y\big),\qquad z=(x,y),
% % \end{equation}
% % and let $g:=\nabla_W\ell(W;z)$ be the per-sample gradient.

% % \paragraph{(1) Hessian.}
% % The Hessian is the exact second derivative of the empirical risk:
% % \begin{equation}
% % H(W):=\nabla_W^2\widehat L_S(W)=\frac{1}{|S|}\sum_{z\in S}\nabla_W^2\ell(W;z).
% % \end{equation}
% % It can be indefinite in general (especially away from minima).

% % \paragraph{(2) Generalized Gauss--Newton (GGN).}
% % For many supervised models, one can write an approximation of the curvature as a Jacobian-sandwich:
% % \begin{equation}
% % G(W)=J(W)^\top\,H_{\text{out}}(W)\,J(W),
% % \end{equation}
% % where $J(W)$ is the Jacobian of the model outputs w.r.t. parameters, and $H_{\text{out}}(W)$ is the Hessian of the loss w.r.t. model outputs.
% % The (generalized) Gauss--Newton matrix is positive semidefinite by construction and coincides with the true Hessian when the model is linear in parameters (or under standard local linearization assumptions).

% % \paragraph{(3) Fisher information (model Fisher).}
% % Assume a probabilistic model $p_W(y\mid x)$ and define the score $s:=\nabla_W\log p_W(y\mid x)$.
% % The (model) Fisher information is
% % \begin{equation}
% % F(W):=\mathbb E_{x\sim\mathcal D_x}\,\mathbb E_{y\sim p_W(\cdot\mid x)}\big[s\,s^\top\big].
% % \end{equation}
% % It is always positive semidefinite and is the curvature underlying the natural gradient~\citep{JMLR:NewInsights}.

% % \paragraph{(4) Empirical Fisher vs expected Fisher.}
% % A commonly used computable surrogate is the \\emph{empirical Fisher}
% % \begin{equation}
% % \widehat F(W):=\frac{1}{|S|}\sum_{(x,y)\in S}\big(\nabla_W\log p_W(y\mid x)\big)\big(\nabla_W\log p_W(y\mid x)\big)^\top.
% % \end{equation}
% % In contrast, the \emph{expected Fisher} is the Fisher $F(W)$ above, where the inner expectation is taken over $y\sim p_W(\cdot\mid x)$ rather than the observed labels.
% % Following~\citep{JMLR:NewInsights}, $\widehat F$ should not be conflated with $F$: they coincide only under restrictive conditions (e.g., when the data labels are actually generated from $p_W(\cdot\mid x)$ and expectations are taken w.r.t. the same distribution).

% % \paragraph{Key relationships (rule of thumb).}
% % \begin{itemize}
% % \item For negative log-likelihood losses, the GGN often matches the (expected) Fisher under standard regularity conditions; both are positive semidefinite.
% % \item The Hessian equals the GGN plus additional terms coming from the curvature of the model mapping; these extra terms vanish for models that are linear in parameters.
% % \item In this note, whenever we write $C_t$, it denotes a task-$t$ second-order curvature proxy at $\widehat W_t$; depending on context it may be instantiated as the Hessian, GGN, or (expected) Fisher, but we keep the symbol $C_t$ fixed for notational uniformity.
% % \end{itemize}

% % \subsection{A unified radius for adversarial and random perturbations}
% % \label{appendix:unified-radius}
% % We will use a single ``radius'' parameter to compare adversarial (ball-constrained) and random (Gaussian) perturbations.

% % \paragraph{Setup (isotropic Gaussian).}
% % Fix a reference point $W\in\mathbb R^d$ and let $\varepsilon\sim\mathcal N(0,\sigma^2 I_d)$.
% % For any confidence level $\alpha\in(0,1)$, define the corresponding high-probability radius
% % \begin{equation}
% % \rho_{\alpha}(\sigma)
% % :=
% % \sigma\,\sqrt{q_{1-\alpha}(\chi^2_d)},
% % \end{equation}
% % where $q_{1-\alpha}(\chi^2_d)$ denotes the $(1-\alpha)$-quantile of the $\chi^2$ distribution with $d$ degrees of freedom.
% % Then $\mathbb P(\|\varepsilon\|\le \rho_{\alpha}(\sigma))\ge 1-\alpha$.
% % We will compare $\mathrm{AS}^{(0)}_S(W,\rho)$ and $\mathrm{RS}^{(0)}_S(W,\sigma^2 I)$ using the same effective radius by setting $\rho=\rho_{\alpha}(\sigma)$, a standard normalization in PAC-Bayes/perturbation-based flatness analyses (see, e.g.,~\citep{germainPACBayesianTheoryMeets2016,tsuzukuNormalizedFlatMinima2020}).

% % \paragraph{Remark (anisotropic Gaussian / subspace).}
% % For $\varepsilon\sim\mathcal N(0,\Sigma)$ supported on a subspace (e.g., LoRA), one may analogously define $\rho_{\alpha}(\Sigma)$ via the quantile of the quadratic form $\varepsilon^\top\Sigma^{\dagger}\varepsilon$ (or equivalently via concentration bounds).

% % \subsection{Proof sketch: why $\mathrm{AS}^{(1)}\ge\mathrm{AS}^{(0)}\ge\mathrm{RS}^{(0)}$}
% % \label{appendix:proof-ordering}
% % We give a two-step argument that makes explicit how the three quantities are compared under a unified radius.

% % \paragraph{Step 1: random sharpness is controlled by adversarial sharpness (same effective radius)}
% % The key observation (used explicitly in SAM-as-Bayes style analyses) is that a supremum over a set upper-bounds the expectation of any random variable supported in that set~\citep{mollenhoffSAMOptimalRelaxation2022}.

% % \begin{lemma}[From ball-supremum to expectation]
% % \label{lemma:sup-to-exp}
% % Let $g:\mathbb R^d\to\mathbb R$ be any function and let $\varepsilon$ be any random vector.
% % Then for any $\rho>0$,
% % \begin{equation}
% % \mathbb E\big[g(\varepsilon)\,\mathbf 1\{\|\varepsilon\|\le\rho\}\big]
% % \le
% % \sup_{\|u\|\le\rho} g(u).
% % \end{equation}
% % \end{lemma}

% % \begin{proof}
% % On the event $\{\|\varepsilon\|\le\rho\}$ we have $g(\varepsilon)\le \sup_{\|u\|\le\rho} g(u)$, and taking expectation yields the claim.
% % \end{proof}

% % Apply Lemma~\ref{lemma:sup-to-exp} to
% % $g(u):=\widehat L_S(W+u)-\widehat L_S(W)$.
% % For the Gaussian perturbation $\varepsilon\sim\mathcal N(0,\sigma^2 I)$ and the unified radius $\rho=\rho_{\alpha}(\sigma)$ from Appendix~\ref{appendix:unified-radius}, we obtain
% % \begin{equation}
% % \mathbb E\big[\widehat L_S(W+\varepsilon)-\widehat L_S(W)\;\mathbf 1\{\|\varepsilon\|\le\rho\}\big]
% % \le
% % \mathrm{AS}^{(0)}_S(W,\rho).
% % \end{equation}
% % If moreover the loss is bounded (e.g., $\ell\in[0,1]$ as assumed in Theorem~\ref{theorem:classical}), then the tail event contributes at most $\alpha$:
% % \begin{equation}
% % \mathrm{RS}^{(0)}_S(W,\sigma^2 I)
% % \le
% % \mathrm{AS}^{(0)}_S(W,\rho_{\alpha}(\sigma))
% % +
% % \alpha.
% % \end{equation}
% % This is the sense in which adversarial zeroth-order sharpness controls the random zeroth-order sharpness when we match them through a single effective radius; compare the discussion in SAM as a relaxation of Bayes~\citep{mollenhoffSAMOptimalRelaxation2022}.

% % \paragraph{Step 2: first-order (gradient) sharpness dominates zeroth-order sharpness}
% % We now relate the loss elevation to the neighborhood gradient norm, following the first-order flatness viewpoint in gradient-norm-aware minimization (GAM)~\citep{Zhang_GAM}.

% % \begin{lemma}[Zeroth-order elevation via first-order sharpness]
% % \label{lemma:as0-vs-as1}
% % Assume $\widehat L_S$ is differentiable.
% % Then for any $\rho>0$,
% % \begin{equation}
% % \mathrm{AS}^{(0)}_S(W,\rho)
% % \le
% % \rho\,\mathrm{AS}^{(1)}_S(W,\rho).
% % \end{equation}
% % \end{lemma}

% % \begin{proof}
% % Fix any $\varepsilon$ with $\|\varepsilon\|\le\rho$.
% % By the fundamental theorem of calculus along the segment $W+s\varepsilon$,
% % \begin{equation}
% % \widehat L_S(W+\varepsilon)-\widehat L_S(W)
% % =
% % \int_0^1 \langle \nabla\widehat L_S(W+s\varepsilon),\varepsilon\rangle\,ds.
% % \end{equation}
% % By Cauchy--Schwarz and $\|\varepsilon\|\le\rho$,
% % \begin{equation}
% % \widehat L_S(W+\varepsilon)-\widehat L_S(W)
% % \le
% % \int_0^1 \|\nabla\widehat L_S(W+s\varepsilon)\|\,\|\varepsilon\|\,ds
% % \le
% % \rho\,\max_{\|u\|\le\rho}\|\nabla\widehat L_S(W+u)\|.
% % \end{equation}
% % Taking the maximum over all $\|\varepsilon\|\le\rho$ yields the result.
% % \end{proof}

% % Combining Lemma~\ref{lemma:as0-vs-as1} with Step 1 (and the same unified radius choice) gives the intended ordering up to a factor $\rho$ and an $\alpha$-tail term:
% % \begin{equation}
% % \mathrm{AS}^{(1)}_S(W,\rho)
% % \;\ge\;
% % \frac{1}{\rho}\,\mathrm{AS}^{(0)}_S(W,\rho)
% % \;\gtrsim\;
% % \frac{1}{\rho}\,\mathrm{RS}^{(0)}_S(W,\sigma^2 I).
% % \end{equation}
% % In the main text we suppress the (typically small) tail term and focus on the qualitative implication that increasingly adversarial objectives impose increasingly strong local control.

% % % \subsection{Proof of Theorem~}
% % % \label{appendix:proof-lora-equivalence}

% % % \begin{proof}
% % % By construction, $\Delta W = PZ = BA$ defines a rank-$r$ matrix. Since both $P$ and $Z$ are trainable without constraints, the function class defined by $PZ$ is identical to that defined by $BA$.
% % % \end{proof}


% % \subsection{Proof of }
% % \label{appendix:construction}

% % \paragraph{Reachability of feature perturbations under admissible PEFT updates.}
% Fix a task $t$ and let $\mathcal W_{\Delta,t}\subseteq\mathbb R^{d_\phi}$ denote the admissible update subspace for the feature extractor parameters $W_\phi$.
% % Let $\Pi_{\Delta,t}$ be the orthogonal projector onto $\mathcal W_{\Delta,t}$, so that
% % \begin{equation}
% % \Pi_{\Delta,t}^2=\Pi_{\Delta,t},\qquad
% % \Pi_{\Delta,t}^\top=\Pi_{\Delta,t},\qquad
% % \mathrm{Im}(\Pi_{\Delta,t})=\mathcal W_{\Delta,t}.
% % \end{equation}

% % Consider the feature representation $h(x)=\phi_{W_\phi}(x)\in\mathbb R^{p}$ and define the Jacobian of the feature extractor
% % \begin{equation}
% % J_\phi(x)
% % :=
% % \frac{\partial \phi_{W_\phi}(x)}{\partial W_\phi}
% % \in\mathbb R^{p\times d_\phi}.
% % \end{equation}
% For a small parameter perturbation $\delta W_\phi$, a first order Taylor expansion of $\phi_{W_\phi}(x)$ around $W_\phi$ gives
% \begin{equation}
% \label{eq:feat_linearization}
% \phi_{W_\phi+\delta W_\phi}(x)
% =
% \phi_{W_\phi}(x)
% +
% J_\phi(x)\,\delta W_\phi
% +
% r(x,\delta W_\phi),
% \qquad
% \|r(x,\delta W_\phi)\|
% =
% o(\|\delta W_\phi\|).
% \end{equation}
% Define the induced feature perturbation
% \begin{equation}
% \delta h(x)
% :=
% \phi_{W_\phi+\delta W_\phi}(x)-\phi_{W_\phi}(x).
% \end{equation}
% Then Eq.~\eqref{eq:feat_linearization} implies
% \begin{equation}
% \label{eq:delta_h_first_order}
% \delta h(x)
% =
% J_\phi(x)\,\delta W_\phi
% +
% r(x,\delta W_\phi)
% \approx
% J_\phi(x)\,\delta W_\phi,
% \end{equation}
% where the approximation holds whenever $\|\delta W_\phi\|$ is sufficiently small so that the remainder is negligible.

% Now assume $\delta W_\phi$ is \emph{admissible}, namely $\delta W_\phi\in\mathcal W_{\Delta,t}$.
% Since $\Pi_{\Delta,t}$ is the projector onto $\mathcal W_{\Delta,t}$, any vector in $\mathcal W_{\Delta,t}$ is invariant under projection:
% \begin{equation}
% \label{eq:proj_invariance}
% \delta W_\phi\in\mathcal W_{\Delta,t}
% \quad\Longrightarrow\quad
% \Pi_{\Delta,t}\delta W_\phi=\delta W_\phi.
% \end{equation}
% Substituting Eq.~\eqref{eq:proj_invariance} into Eq.~\eqref{eq:delta_h_first_order} yields the equivalent first order expression
% \begin{equation}
% \label{eq:delta_h_with_proj}
% \delta h(x)
% \approx
% J_\phi(x)\,\Pi_{\Delta,t}\delta W_\phi.
% \end{equation}

% Finally, define the \emph{reachable feature change subspace} at input $x$ by
% \begin{equation}
% \label{eq:reachable_def_recap}
% \mathcal R_t(x)
% :=
% \mathrm{Im}\!\bigl(J_\phi(x)\Pi_{\Delta,t}\bigr)
% =
% \Bigl\{\,J_\phi(x)\Pi_{\Delta,t}v\ \big|\ v\in\mathbb R^{d_\phi}\Bigr\}
% \subseteq \mathbb R^{p}.
% \end{equation}
% By the set definition in Eq.~\eqref{eq:reachable_def_recap}, any vector of the form $J_\phi(x)\Pi_{\Delta,t}v$ belongs to $\mathcal R_t(x)$.
% Applying this with $v=\delta W_\phi$ and using Eq.~\eqref{eq:delta_h_with_proj} proves that, for any admissible $\delta W_\phi\in\mathcal W_{\Delta,t}$,
% \begin{equation}
% \delta h(x)\approx J_\phi(x)\Pi_{\Delta,t}\delta W_\phi \in \mathcal R_t(x).
% \end{equation}



% % % \section{discussion on weight distribution drift}
% % \section{Information geometric decomposition for comparing two solutions}

% % Let $W_0,W_1$ be two parameter states (matrices or vectors). We compare them through a local Gaussian information model on the vectorized parameter
% % \begin{equation}
% % w_i=\mathrm{vec}(W_i)\in\mathbb R^D,\qquad p_i(w)=\mathcal N(w_i,\Sigma_i),\qquad i\in{0,1}.
% % \end{equation}
% % The information discrepancy is quantified by the Gaussian KL divergence, which admits the standard decomposition
% % \begin{equation}
% % \mathrm{KL}(p_1\Vert p_0)
% % =
% % \frac12\Big(
% % \underbrace{(w_1-w_0)^\top\Sigma_0^{-1}(w_1-w_0)}_{\triangleq\ 2D_{\mathrm{mean}}(W_1,W_0)}
% % +
% % \underbrace{\operatorname{tr}(\Sigma_0^{-1}\Sigma_1)-D-\log\det(\Sigma_0^{-1}\Sigma_1)}_{\triangleq\ 2D_{\mathrm{cov}}(W_1,W_0)}
% % \Big).
% % \end{equation}
% % Here $D_{\mathrm{mean}}$ is a whitened displacement (a Mahalanobis squared distance), while $D_{\mathrm{cov}}$ measures the change in local uncertainty geometry (shape and orientation) between $\Sigma_0$ and $\Sigma_1$. 

% % \subsection{Why isotropic covariance cannot express subspace mismatch}

% % If $\Sigma_0=\sigma^2 I$ is isotropic, then
% % \begin{equation}
% % D_{\mathrm{mean}}(W_1,W_0)=\frac{1}{2\sigma^2}\lVert w_1-w_0\rVert_2^2,
% % \end{equation}
% % and the covariance term reduces to a spectral functional of $\Sigma_1$ (it depends on eigenvalues of $\Sigma_1$ but is invariant to rotations relative to $W_0$). Consequently, isotropic modeling cannot attribute any part of the complexity to a “directional mismatch” between two structured subspaces. 

% % \subsection{Low rank structured covariance and Grassmannian drift}

% % To make the covariance drift sensitive to directional structure, restrict covariances to an isotropic plus subspace energy form
% % \begin{equation}
% % \Sigma_i=\sigma^2 I+\tau^2 P_i,\qquad P_i=Y_iY_i^\top,\qquad Y_i\in\mathbb R^{D\times m},\quad Y_i^\top Y_i=I_m.
% % \end{equation}
% % The relative position of $P_0$ and $P_1$ is characterized by the principal angles ${\theta_k}_{k=1}^m$ between $\mathrm{span}(Y_0)$ and $\mathrm{span}(Y_1)$. A canonical, directly computable Grassmannian proxy is the projection metric
% % \begin{equation}
% % D_{\mathrm{scope}}(W_1,W_0)\triangleq d_P^2(Y_0,Y_1)
% % =
% %  m-\lVert Y_0^\top Y_1\rVert_F^2
% % \sum_{k=1}^m\sin^2\theta_k.
% % \end{equation}
% % Under the structured covariance assumption above, the log determinant ratio cancels and the covariance drift depends on the subspace mismatch only through ${\theta_k}$, yielding (up to constants depending on $\sigma,\tau$)
% % \begin{equation}
% % D_{\mathrm{cov}}(W_1,W_0)
% % =
% % \text{constant}
% % +
% % c(\sigma,\tau),d_P^2(Y_0,Y_1),
% % \qquad c(\sigma,\tau)>0,
% % \end{equation}
% % so $D_{\mathrm{scope}}$ is not an ad hoc diagnostic but an explicit geometric surrogate for the direction dependent component of $D_{\mathrm{cov}}$.

% % \subsection{Connecting “prior information injection” to the mean drift geometry}

% % Injecting past task sensitivity changes the geometry of the displacement term by replacing an isotropic precision with a structured precision. A generic parameterization is
% % \begin{equation}
% % \Sigma_P^{-1}=\lambda I+\alpha G,
% % \end{equation}
% % where $G\succeq 0$ encodes protected directions, for example a Fisher based sensitivity operator. Then the mean drift becomes a sensitivity weighted displacement
% % \begin{equation}
% % D_{\mathrm{mean}}=\frac12(w_1-\mu_P)^\top(\lambda I+\alpha G)(w_1-\mu_P),
% % \end{equation}
% % which discourages movement along high sensitivity directions while leaving orthogonal directions available for adaptation. This is the precise sense in which “protecting important weights” is equivalent to defining a protected metric or protected subspace rather than locking individual coordinates.

% % \subsection{Operationalizing $Y_t$ in PEFT and continual learning}

% % The above quantities become meaningful only after specifying how $Y_t$ is constructed. Viable choices that remain aligned with PEFT and continual learning interpretations include:

% % 1. Curvature principal subspace: $Y_t$ as the top $m$ eigenvectors of a task curvature proxy (Hessian or Fisher), optionally restricted to the trainable adapter coordinates.
% % 2. Task permissible update geometry: $Y_t$ as an orthonormal basis of the task admissible adapter update space $\mathcal W_{\Delta,t}$, or a principal basis derived from the LoRA induced tangent geometry at the final adapter factors.
% % 3. Noise geometry: $Y_t$ as the top subspace of an estimated optimizer noise covariance.

% % With a fixed convention for $Y_t$, the triplet
% % \begin{equation}
% % \mathrm{KL}(p_t\Vert p_{t-1})=D_{\mathrm{mean}}(W_t,W_{t-1})+D_{\mathrm{cov}}(W_t,W_{t-1}),
% % \qquad
% % D_{\mathrm{scope}}(W_t,W_{t-1})=d_P^2(Y_{t-1},Y_t),
% % \end{equation}
% % provides a compact, reproducible summary of how much information is written (mean drift) and how the dominant geometric directions rotate across tasks (scope drift).


% % % \subsection{Proof of Proposition~\ref{prop:non-equiv}}
% % % \label{appendix:proof-non-equivalence}

% % % \begin{proof}
% % % Assume a fixed $A$ exists such that $BA = PZ$ holds for all $Z \sim \mathcal{N}(0, \sigma^2 I)$. Then $BA$ is deterministic, while $PZ$ is stochastic over the support of $Z$, leading to a contradiction almost surely. Therefore, no such $A$ exists.
% % % \end{proof}





% % \section{Technical Appendices and Supplementary Material}
% % Technical appendices with additional results, figures, graphs and proofs may be submitted with the paper submission before the full submission deadline (see above), or as a separate PDF in the ZIP file below before the supplementary material deadline. There is no page limit for the technical appendices.

% % %%%%%%%%%%%%%%%%%%%%%%%%%%%%%%%%%%%%%%%%%%%%%%%%%%%%%%%%%%%%

% % %\newpage
% % % \section*{NeurIPS Paper Checklist}

% % % %%% BEGIN INSTRUCTIONS %%%
% % % The checklist is designed to encourage best practices for responsible machine learning research, addressing issues of reproducibility, transparency, research ethics, and societal impact. Do not remove the checklist: {\bf The papers not including the checklist will be desk rejected.} The checklist should follow the references and follow the (optional) supplemental material.  The checklist does NOT count towards the page
% % % limit. 

% % % Please read the checklist guidelines carefully for information on how to answer these questions. For each question in the checklist:
% % % \begin{itemize}
% % %     \item You should answer \answerYes{}, \answerNo{}, or \answerNA{}.
% % %     \item \answerNA{} means either that the question is Not Applicable for that particular paper or the relevant information is Not Available.
% % %     \item Please provide a short (1–2 sentence) justification right after your answer (even for NA). 
% % %    % \item {\bf The papers not including the checklist will be desk rejected.}
% % % \end{itemize}

% % % {\bf The checklist answers are an integral part of your paper submission.} They are visible to the reviewers, area chairs, senior area chairs, and ethics reviewers. You will be asked to also include it (after eventual revisions) with the final version of your paper, and its final version will be published with the paper.

% % % The reviewers of your paper will be asked to use the checklist as one of the factors in their evaluation. While "\answerYes{}" is generally preferable to "\answerNo{}", it is perfectly acceptable to answer "\answerNo{}" provided a proper justification is given (e.g., "error bars are not reported because it would be too computationally expensive" or "we were unable to find the license for the dataset we used"). In general, answering "\answerNo{}" or "\answerNA{}" is not grounds for rejection. While the questions are phrased in a binary way, we acknowledge that the true answer is often more nuanced, so please just use your best judgment and write a justification to elaborate. All supporting evidence can appear either in the main paper or the supplemental material, provided in appendix. If you answer \answerYes{} to a question, in the justification please point to the section(s) where related material for the question can be found.

% % % IMPORTANT, please:
% % % \begin{itemize}
% % %     \item {\bf Delete this instruction block, but keep the section heading ``NeurIPS Paper Checklist"},
% % %     \item  {\bf Keep the checklist subsection headings, questions/answers and guidelines below.}
% % %     \item {\bf Do not modify the questions and only use the provided macros for your answers}.
% % % \end{itemize} 
 

% % % %%% END INSTRUCTIONS %%%


% % % \begin{enumerate}

% % % \item {\bf Claims}
% % %     \item[] Question: Do the main claims made in the abstract and introduction accurately reflect the paper's contributions and scope?
% % %     \item[] Answer: \answerTODO{} % Replace by \answerYes{}, \answerNo{}, or \answerNA{}.
% % %     \item[] Justification: \justificationTODO{}
% % %     \item[] Guidelines:
% % %     \begin{itemize}
% % %         \item The answer NA means that the abstract and introduction do not include the claims made in the paper.
% % %         \item The abstract and/or introduction should clearly state the claims made, including the contributions made in the paper and important assumptions and limitations. A No or NA answer to this question will not be perceived well by the reviewers. 
% % %         \item The claims made should match theoretical and experimental results, and reflect how much the results can be expected to generalize to other settings. 
% % %         \item It is fine to include aspirational goals as motivation as long as it is clear that these goals are not attained by the paper. 
% % %     \end{itemize}

% % % \item {\bf Limitations}
% % %     \item[] Question: Does the paper discuss the limitations of the work performed by the authors?
% % %     \item[] Answer: \answerTODO{} % Replace by \answerYes{}, \answerNo{}, or \answerNA{}.
% % %     \item[] Justification: \justificationTODO{}
% % %     \item[] Guidelines:
% % %     \begin{itemize}
% % %         \item The answer NA means that the paper has no limitation while the answer No means that the paper has limitations, but those are not discussed in the paper. 
% % %         \item The authors are encouraged to create a separate "Limitations" section in their paper.
% % %         \item The paper should point out any strong assumptions and how robust the results are to violations of these assumptions (e.g., independence assumptions, noiseless settings, model well-specification, asymptotic approximations only holding locally). The authors should reflect on how these assumptions might be violated in practice and what the implications would be.
% % %         \item The authors should reflect on the scope of the claims made, e.g., if the approach was only tested on a few datasets or with a few runs. In general, empirical results often depend on implicit assumptions, which should be articulated.
% % %         \item The authors should reflect on the factors that influence the performance of the approach. For example, a facial recognition algorithm may perform poorly when image resolution is low or images are taken in low lighting. Or a speech-to-text system might not be used reliably to provide closed captions for online lectures because it fails to handle technical jargon.
% % %         \item The authors should discuss the computational efficiency of the proposed algorithms and how they scale with dataset size.
% % %         \item If applicable, the authors should discuss possible limitations of their approach to address problems of privacy and fairness.
% % %         \item While the authors might fear that complete honesty about limitations might be used by reviewers as grounds for rejection, a worse outcome might be that reviewers discover limitations that aren't acknowledged in the paper. The authors should use their best judgment and recognize that individual actions in favor of transparency play an important role in developing norms that preserve the integrity of the community. Reviewers will be specifically instructed to not penalize honesty concerning limitations.
% % %     \end{itemize}

% % % \item {\bf Theory assumptions and proofs}
% % %     \item[] Question: For each theoretical result, does the paper provide the full set of assumptions and a complete (and correct) proof?
% % %     \item[] Answer: \answerTODO{} % Replace by \answerYes{}, \answerNo{}, or \answerNA{}.
% % %     \item[] Justification: \justificationTODO{}
% % %     \item[] Guidelines:
% % %     \begin{itemize}
% % %         \item The answer NA means that the paper does not include theoretical results. 
% % %         \item All the theorems, formulas, and proofs in the paper should be numbered and cross-referenced.
% % %         \item All assumptions should be clearly stated or referenced in the statement of any theorems.
% % %         \item The proofs can either appear in the main paper or the supplemental material, but if they appear in the supplemental material, the authors are encouraged to provide a short proof sketch to provide intuition. 
% % %         \item Inversely, any informal proof provided in the core of the paper should be complemented by formal proofs provided in appendix or supplemental material.
% % %         \item Theorems and Lemmas that the proof relies upon should be properly referenced. 
% % %     \end{itemize}

% % %     \item {\bf Experimental result reproducibility}
% % %     \item[] Question: Does the paper fully disclose all the information needed to reproduce the main experimental results of the paper to the extent that it affects the main claims and/or conclusions of the paper (regardless of whether the code and data are provided or not)?
% % %     \item[] Answer: \answerTODO{} % Replace by \answerYes{}, \answerNo{}, or \answerNA{}.
% % %     \item[] Justification: \justificationTODO{}
% % %     \item[] Guidelines:
% % %     \begin{itemize}
% % %         \item The answer NA means that the paper does not include experiments.
% % %         \item If the paper includes experiments, a No answer to this question will not be perceived well by the reviewers: Making the paper reproducible is important, regardless of whether the code and data are provided or not.
% % %         \item If the contribution is a dataset and/or model, the authors should describe the steps taken to make their results reproducible or verifiable. 
% % %         \item Depending on the contribution, reproducibility can be accomplished in various ways. For example, if the contribution is a novel architecture, describing the architecture fully might suffice, or if the contribution is a specific model and empirical evaluation, it may be necessary to either make it possible for others to replicate the model with the same dataset, or provide access to the model. In general. releasing code and data is often one good way to accomplish this, but reproducibility can also be provided via detailed instructions for how to replicate the results, access to a hosted model (e.g., in the case of a large language model), releasing of a model checkpoint, or other means that are appropriate to the research performed.
% % %         \item While NeurIPS does not require releasing code, the conference does require all submissions to provide some reasonable avenue for reproducibility, which may depend on the nature of the contribution. For example
% % %         \begin{enumerate}
% % %             \item If the contribution is primarily a new algorithm, the paper should make it clear how to reproduce that algorithm.
% % %             \item If the contribution is primarily a new model architecture, the paper should describe the architecture clearly and fully.
% % %             \item If the contribution is a new model (e.g., a large language model), then there should either be a way to access this model for reproducing the results or a way to reproduce the model (e.g., with an open-source dataset or instructions for how to construct the dataset).
% % %             \item We recognize that reproducibility may be tricky in some cases, in which case authors are welcome to describe the particular way they provide for reproducibility. In the case of closed-source models, it may be that access to the model is limited in some way (e.g., to registered users), but it should be possible for other researchers to have some path to reproducing or verifying the results.
% % %         \end{enumerate}
% % %     \end{itemize}


% % % \item {\bf Open access to data and code}
% % %     \item[] Question: Does the paper provide open access to the data and code, with sufficient instructions to faithfully reproduce the main experimental results, as described in supplemental material?
% % %     \item[] Answer: \answerTODO{} % Replace by \answerYes{}, \answerNo{}, or \answerNA{}.
% % %     \item[] Justification: \justificationTODO{}
% % %     \item[] Guidelines:
% % %     \begin{itemize}
% % %         \item The answer NA means that paper does not include experiments requiring code.
% % %         \item Please see the NeurIPS code and data submission guidelines (\url{https://nips.cc/public/guides/CodeSubmissionPolicy}) for more details.
% % %         \item While we encourage the release of code and data, we understand that this might not be possible, so “No” is an acceptable answer. Papers cannot be rejected simply for not including code, unless this is central to the contribution (e.g., for a new open-source benchmark).
% % %         \item The instructions should contain the exact command and environment needed to run to reproduce the results. See the NeurIPS code and data submission guidelines (\url{https://nips.cc/public/guides/CodeSubmissionPolicy}) for more details.
% % %         \item The authors should provide instructions on data access and preparation, including how to access the raw data, preprocessed data, intermediate data, and generated data, etc.
% % %         \item The authors should provide scripts to reproduce all experimental results for the new proposed method and baselines. If only a subset of experiments are reproducible, they should state which ones are omitted from the script and why.
% % %         \item At submission time, to preserve anonymity, the authors should release anonymized versions (if applicable).
% % %         \item Providing as much information as possible in supplemental material (appended to the paper) is recommended, but including URLs to data and code is permitted.
% % %     \end{itemize}


% % % \item {\bf Experimental setting/details}
% % %     \item[] Question: Does the paper specify all the training and test details (e.g., data splits, hyperparameters, how they were chosen, type of optimizer, etc.) necessary to understand the results?
% % %     \item[] Answer: \answerTODO{} % Replace by \answerYes{}, \answerNo{}, or \answerNA{}.
% % %     \item[] Justification: \justificationTODO{}
% % %     \item[] Guidelines:
% % %     \begin{itemize}
% % %         \item The answer NA means that the paper does not include experiments.
% % %         \item The experimental setting should be presented in the core of the paper to a level of detail that is necessary to appreciate the results and make sense of them.
% % %         \item The full details can be provided either with the code, in appendix, or as supplemental material.
% % %     \end{itemize}

% % % \item {\bf Experiment statistical significance}
% % %     \item[] Question: Does the paper report error bars suitably and correctly defined or other appropriate information about the statistical significance of the experiments?
% % %     \item[] Answer: \answerTODO{} % Replace by \answerYes{}, \answerNo{}, or \answerNA{}.
% % %     \item[] Justification: \justificationTODO{}
% % %     \item[] Guidelines:
% % %     \begin{itemize}
% % %         \item The answer NA means that the paper does not include experiments.
% % %         \item The authors should answer "Yes" if the results are accompanied by error bars, confidence intervals, or statistical significance tests, at least for the experiments that support the main claims of the paper.
% % %         \item The factors of variability that the error bars are capturing should be clearly stated (for example, train/test split, initialization, random drawing of some parameter, or overall run with given experimental conditions).
% % %         \item The method for calculating the error bars should be explained (closed form formula, call to a library function, bootstrap, etc.)
% % %         \item The assumptions made should be given (e.g., Normally distributed errors).
% % %         \item It should be clear whether the error bar is the standard deviation or the standard error of the mean.
% % %         \item It is OK to report 1-sigma error bars, but one should state it. The authors should preferably report a 2-sigma error bar than state that they have a 96\% CI, if the hypothesis of Normality of errors is not verified.
% % %         \item For asymmetric distributions, the authors should be careful not to show in tables or figures symmetric error bars that would yield results that are out of range (e.g. negative error rates).
% % %         \item If error bars are reported in tables or plots, The authors should explain in the text how they were calculated and reference the corresponding figures or tables in the text.
% % %     \end{itemize}

% % % \item {\bf Experiments compute resources}
% % %     \item[] Question: For each experiment, does the paper provide sufficient information on the computer resources (type of compute workers, memory, time of execution) needed to reproduce the experiments?
% % %     \item[] Answer: \answerTODO{} % Replace by \answerYes{}, \answerNo{}, or \answerNA{}.
% % %     \item[] Justification: \justificationTODO{}
% % %     \item[] Guidelines:
% % %     \begin{itemize}
% % %         \item The answer NA means that the paper does not include experiments.
% % %         \item The paper should indicate the type of compute workers CPU or GPU, internal cluster, or cloud provider, including relevant memory and storage.
% % %         \item The paper should provide the amount of compute required for each of the individual experimental runs as well as estimate the total compute. 
% % %         \item The paper should disclose whether the full research project required more compute than the experiments reported in the paper (e.g., preliminary or failed experiments that didn't make it into the paper). 
% % %     \end{itemize}
    
% % % \item {\bf Code of ethics}
% % %     \item[] Question: Does the research conducted in the paper conform, in every respect, with the NeurIPS Code of Ethics \url{https://neurips.cc/public/EthicsGuidelines}?
% % %     \item[] Answer: \answerTODO{} % Replace by \answerYes{}, \answerNo{}, or \answerNA{}.
% % %     \item[] Justification: \justificationTODO{}
% % %     \item[] Guidelines:
% % %     \begin{itemize}
% % %         \item The answer NA means that the authors have not reviewed the NeurIPS Code of Ethics.
% % %         \item If the authors answer No, they should explain the special circumstances that require a deviation from the Code of Ethics.
% % %         \item The authors should make sure to preserve anonymity (e.g., if there is a special consideration due to laws or regulations in their jurisdiction).
% % %     \end{itemize}


% % % \item {\bf Broader impacts}
% % %     \item[] Question: Does the paper discuss both potential positive societal impacts and negative societal impacts of the work performed?
% % %     \item[] Answer: \answerTODO{} % Replace by \answerYes{}, \answerNo{}, or \answerNA{}.
% % %     \item[] Justification: \justificationTODO{}
% % %     \item[] Guidelines:
% % %     \begin{itemize}
% % %         \item The answer NA means that there is no societal impact of the work performed.
% % %         \item If the authors answer NA or No, they should explain why their work has no societal impact or why the paper does not address societal impact.
% % %         \item Examples of negative societal impacts include potential malicious or unintended uses (e.g., disinformation, generating fake profiles, surveillance), fairness considerations (e.g., deployment of technologies that could make decisions that unfairly impact specific groups), privacy considerations, and security considerations.
% % %         \item The conference expects that many papers will be foundational research and not tied to particular applications, let alone deployments. However, if there is a direct path to any negative applications, the authors should point it out. For example, it is legitimate to point out that an improvement in the quality of generative models could be used to generate deepfakes for disinformation. On the other hand, it is not needed to point out that a generic algorithm for optimizing neural networks could enable people to train models that generate Deepfakes faster.
% % %         \item The authors should consider possible harms that could arise when the technology is being used as intended and functioning correctly, harms that could arise when the technology is being used as intended but gives incorrect results, and harms following from (intentional or unintentional) misuse of the technology.
% % %         \item If there are negative societal impacts, the authors could also discuss possible mitigation strategies (e.g., gated release of models, providing defenses in addition to attacks, mechanisms for monitoring misuse, mechanisms to monitor how a system learns from feedback over time, improving the efficiency and accessibility of ML).
% % %     \end{itemize}
    
% % % \item {\bf Safeguards}
% % %     \item[] Question: Does the paper describe safeguards that have been put in place for responsible release of data or models that have a high risk for misuse (e.g., pretrained language models, image generators, or scraped datasets)?
% % %     \item[] Answer: \answerTODO{} % Replace by \answerYes{}, \answerNo{}, or \answerNA{}.
% % %     \item[] Justification: \justificationTODO{}
% % %     \item[] Guidelines:
% % %     \begin{itemize}
% % %         \item The answer NA means that the paper poses no such risks.
% % %         \item Released models that have a high risk for misuse or dual-use should be released with necessary safeguards to allow for controlled use of the model, for example by requiring that users adhere to usage guidelines or restrictions to access the model or implementing safety filters. 
% % %         \item Datasets that have been scraped from the Internet could pose safety risks. The authors should describe how they avoided releasing unsafe images.
% % %         \item We recognize that providing effective safeguards is challenging, and many papers do not require this, but we encourage authors to take this into account and make a best faith effort.
% % %     \end{itemize}

% % % \item {\bf Licenses for existing assets}
% % %     \item[] Question: Are the creators or original owners of assets (e.g., code, data, models), used in the paper, properly credited and are the license and terms of use explicitly mentioned and properly respected?
% % %     \item[] Answer: \answerTODO{} % Replace by \answerYes{}, \answerNo{}, or \answerNA{}.
% % %     \item[] Justification: \justificationTODO{}
% % %     \item[] Guidelines:
% % %     \begin{itemize}
% % %         \item The answer NA means that the paper does not use existing assets.
% % %         \item The authors should cite the original paper that produced the code package or dataset.
% % %         \item The authors should state which version of the asset is used and, if possible, include a URL.
% % %         \item The name of the license (e.g., CC-BY 4.0) should be included for each asset.
% % %         \item For scraped data from a particular source (e.g., website), the copyright and terms of service of that source should be provided.
% % %         \item If assets are released, the license, copyright information, and terms of use in the package should be provided. For popular datasets, \url{paperswithcode.com/datasets} has curated licenses for some datasets. Their licensing guide can help determine the license of a dataset.
% % %         \item For existing datasets that are re-packaged, both the original license and the license of the derived asset (if it has changed) should be provided.
% % %         \item If this information is not available online, the authors are encouraged to reach out to the asset's creators.
% % %     \end{itemize}

% % % \item {\bf New assets}
% % %     \item[] Question: Are new assets introduced in the paper well documented and is the documentation provided alongside the assets?
% % %     \item[] Answer: \answerTODO{} % Replace by \answerYes{}, \answerNo{}, or \answerNA{}.
% % %     \item[] Justification: \justificationTODO{}
% % %     \item[] Guidelines:
% % %     \begin{itemize}
% % %         \item The answer NA means that the paper does not release new assets.
% % %         \item Researchers should communicate the details of the dataset/code/model as part of their submissions via structured templates. This includes details about training, license, limitations, etc. 
% % %         \item The paper should discuss whether and how consent was obtained from people whose asset is used.
% % %         \item At submission time, remember to anonymize your assets (if applicable). You can either create an anonymized URL or include an anonymized zip file.
% % %     \end{itemize}

% % % \item {\bf Crowdsourcing and research with human subjects}
% % %     \item[] Question: For crowdsourcing experiments and research with human subjects, does the paper include the full text of instructions given to participants and screenshots, if applicable, as well as details about compensation (if any)? 
% % %     \item[] Answer: \answerTODO{} % Replace by \answerYes{}, \answerNo{}, or \answerNA{}.
% % %     \item[] Justification: \justificationTODO{}
% % %     \item[] Guidelines:
% % %     \begin{itemize}
% % %         \item The answer NA means that the paper does not involve crowdsourcing nor research with human subjects.
% % %         \item Including this information in the supplemental material is fine, but if the main contribution of the paper involves human subjects, then as much detail as possible should be included in the main paper. 
% % %         \item According to the NeurIPS Code of Ethics, workers involved in data collection, curation, or other labor should be paid at least the minimum wage in the country of the data collector. 
% % %     \end{itemize}

% % % \item {\bf Institutional review board (IRB) approvals or equivalent for research with human subjects}
% % %     \item[] Question: Does the paper describe potential risks incurred by study participants, whether such risks were disclosed to the subjects, and whether Institutional Review Board (IRB) approvals (or an equivalent approval/review based on the requirements of your country or institution) were obtained?
% % %     \item[] Answer: \answerTODO{} % Replace by \answerYes{}, \answerNo{}, or \answerNA{}.
% % %     \item[] Justification: \justificationTODO{}
% % %     \item[] Guidelines:
% % %     \begin{itemize}
% % %         \item The answer NA means that the paper does not involve crowdsourcing nor research with human subjects.
% % %         \item Depending on the country in which research is conducted, IRB approval (or equivalent) may be required for any human subjects research. If you obtained IRB approval, you should clearly state this in the paper. 
% % %         \item We recognize that the procedures for this may vary significantly between institutions and locations, and we expect authors to adhere to the NeurIPS Code of Ethics and the guidelines for their institution. 
% % %         \item For initial submissions, do not include any information that would break anonymity (if applicable), such as the institution conducting the review.
% % %     \end{itemize}

% % % \item {\bf Declaration of LLM usage}
% % %     \item[] Question: Does the paper describe the usage of LLMs if it is an important, original, or non-standard component of the core methods in this research? Note that if the LLM is used only for writing, editing, or formatting purposes and does not impact the core methodology, scientific rigorousness, or originality of the research, declaration is not required.
% % %     %this research? 
% % %     \item[] Answer: \answerTODO{} % Replace by \answerYes{}, \answerNo{}, or \answerNA{}.
% % %     \item[] Justification: \justificationTODO{}
% % %     \item[] Guidelines:
% % %     \begin{itemize}
% % %         \item The answer NA means that the core method development in this research does not involve LLMs as any important, original, or non-standard components.
% % %         \item Please refer to our LLM policy (\url{https://neurips.cc/Conferences/2025/LLM}) for what should or should not be described.
% % %     \end{itemize}

% % % \end{enumerate}

\section{Appendix}

\section{Proof of Proposition X}

We consider the constrained optimization problem

\begin{equation}
\min_{\Sigma_t}
\operatorname{tr}(F_{1,t}\Sigma_t)
-
\beta
\operatorname{tr}(F_{2,t}\Sigma_t)
\end{equation}

subject to

\begin{equation}
\operatorname{tr}(\Sigma_t)=\sigma^2,
\qquad
\Sigma_t \succeq 0 .
\end{equation}

Define

\begin{equation}
M_t = F_{1,t}-\beta F_{2,t}.
\end{equation}

The objective becomes

\begin{equation}
\min_{\Sigma_t} \operatorname{tr}(M_t \Sigma_t).
\end{equation}

Introducing the Lagrangian

\begin{equation}
\mathcal{L}(\Sigma_t,\lambda)
=
\operatorname{tr}(M_t \Sigma_t)
+
\lambda(\operatorname{tr}(\Sigma_t)-\sigma^2).
\end{equation}

Taking the derivative with respect to $\Sigma_t$ and setting it to zero gives

\begin{equation}
M_t + \lambda I = 0.
\end{equation}

To ensure a valid covariance matrix and stable solution, we introduce a regularization term $\epsilon I$ and obtain

\begin{equation}
\Sigma_t^\star
\propto
(M_t + \epsilon I)^{-1}.
\end{equation}

Finally, enforcing the energy constraint $\operatorname{tr}(\Sigma_t)=\sigma^2$ yields

\begin{equation}
\Sigma_t^\star
=
\sigma^2
(M_t + \epsilon I)^{-1}.
\end{equation}

This completes the proof.



% \paragraph{Vector-to-matrix formulation of the local quadratic model.} 

% Let the parameter increment be represented in matrix form 
% \(\Delta W\), and define the corresponding vectorized displacement 
% \begin{equation} w = \theta_1 + \mathrm{vec}(\Delta W), \qquad \delta ;:=; w-\theta_1 = \mathrm{vec}(\Delta W). \label{eq:vec_matrix_bridge} 
% \end{equation} 
% Consider the empirical risk \(\hat L_{S_1}(w)\) as a scalar function of \(w\in\mathbb R^p\). Its second-order Taylor expansion at \(w=\theta_1\) yields 
% \begin{equation} \hat L_{S_1}(\theta_1+\delta)
% =
% \hat L_{S_1}(\theta_1)
% +
% \nabla_w \hat L_{S_1}(\theta_1)^\top \delta
% +
% \frac12,\delta^\top \nabla_w^2 \hat L_{S_1}(\theta_1),\delta
% +
% R_3,
% \label{eq:taylor_delta_form}
% \end{equation}
% where \(R_3\) denotes the third-order remainder term.

% Now substitute \(\delta=\mathrm{vec}(\Delta W)\). Define the gradient and Hessian with respect to the matrix increment \(\Delta W\) at \(\Delta W=0\) by 
% \begin{equation} g_{1,S_1} := \left.\nabla_{\Delta W}\hat L_{S_1}(\theta_1+\Delta W)\right|_{\Delta W=0}, \qquad C_{1,S_1}^{\Delta} := \left.\nabla_{\Delta W}^{2}\hat L_{S_1}(\theta_1+\Delta W)\right|_{\Delta W=0}. \label{eq:grad_hess_DeltaW} 
% \end{equation} 
% Using the standard identity between vectorization and the Frobenius inner product, \begin{equation} \langle A,B\rangle_F
% \mathrm{vec}(A)^\top \mathrm{vec}(B), \label{eq:frob_vec_identity} \end{equation} the linear term in \eqref{eq:taylor_delta_form} can be written equivalently as \begin{equation} \nabla_w \hat L_{S_1}(\theta_1)^\top= \mathrm{vec}(\Delta W)
% \langle g_{1,S_1},\Delta W\rangle_F. \label{eq:first_order_frob} \end{equation} Similarly, the quadratic term admits the matrix-form representation 
% \begin{equation} \frac12,\mathrm{vec}(\Delta W)^\top C_{1,S_1}^{\Delta},\mathrm{vec}(\Delta W), \label{eq:second_order_vec} \end{equation} 
% which depends only on the symmetric part of \(C_{1,S_1}^{\Delta}\) since \(\delta^\top C \delta=\delta^\top \tfrac12(C+C^\top)\delta\). 
% Hence we define the symmetric Hessian
% \begin{equation} \widetilde C_{1,S_1}^{\Delta} := \frac12\left(C_{1,S_1}^{\Delta}+{C_{1,S_1}^{\Delta}}^\top\right), \label{eq:sym_hess_def} \end{equation} 
% and obtain the matrix-form second-order expansion \begin{equation} \hat L_{S_1}(\theta_1+\Delta W)-\hat L_{S_1}(\theta_1)=
% \langle g_{1,S_1},\Delta W\rangle_F
% +
% \frac12,\mathrm{vec}(\Delta W)^\top \widetilde C_{1,S_1}^{\Delta},\mathrm{vec}(\Delta W)
% +
% R_3.
% \label{eq:taylor_matrix_form_final}
% \end{equation}
% If the Hessian is \(L_H\)-Lipschitz in a neighborhood of \(\Delta W=0\) under the chosen norm, then the remainder satisfies the standard bound
% \begin{equation}
% |R_3|
% \le
% \frac{L_H}{6}|\Delta W|_F^3.
% \label{eq:remainder_bound_final}
% \end{equation}



% \paragraph{Part II: From point update to distributional update.}

% We now lift the local quadratic model from a pointwise parameter update to a distributional update.

% Let
% \begin{equation}
% \mathcal P_1=\mathcal N(\theta_1,\Sigma_1),
% \qquad
% \mathcal P_2=\mathcal N(\theta_2,\Sigma_2),
% \qquad
% \delta_{1\to2}:=\theta_2-\theta_1=\mathrm{vec}(\Delta W_2).
% \label{eq:dist_setup}
% \end{equation}

% Define the distributional forgetting on the old task \(S_1\) as
% \begin{equation}
% \Delta^{\mathrm{dist}}_{1\to2}
% :=
% \mathbb E_{w\sim\mathcal P_2}\hat L_{S_1}(w)
% -\mathbb E_{w\sim\mathcal P_1}\hat L_{S_1}(w).
% \label{eq:dist_forgetting_def}
% \end{equation}

% Step 1: Apply the same local quadratic model

% Using the second-order expansion at \(w=\theta_1\) derived in Part I,
% \begin{equation}
% \hat L_{S_1}(w)
% =
% \hat L_{S_1}(\theta_1)
% +
% g_{1,S_1}^{\top}(w-\theta_1)
% +
% \frac12 (w-\theta_1)^{\top}\widetilde C_{1,S_1}^{\Delta}(w-\theta_1)
% +
% R_3(w),
% \label{eq:local_quad_model}
% \end{equation}
% where \(g_{1,S_1}\) and \(\widetilde C_{1,S_1}^{\Delta}\) are defined in Part I.

% Step 2: Take expectation under a Gaussian parameter distribution

% For a Gaussian random vector
% $
% w\sim\mathcal N(\theta,\Sigma),
% $
% we use the standard identities:
% \begin{equation}
% \mathbb E[w-\theta_1]=\theta-\theta_1,
% \label{eq:first_moment}
% \end{equation}
% \begin{equation}
% \mathbb E\big[(w-\theta_1)(w-\theta_1)^\top\big]
% =
% (\theta-\theta_1)(\theta-\theta_1)^\top
% +
% \Sigma.
% \label{eq:second_moment}
% \end{equation}

% Substituting \eqref{eq:first_moment}–\eqref{eq:second_moment} into \eqref{eq:local_quad_model} yields
% \begin{align}
% \mathbb E_{w\sim\mathcal N(\theta,\Sigma)}\hat L_{S_1}(w)
% &=
% \hat L_{S_1}(\theta_1)
% +
% g_{1,S_1}^{\top}(\theta-\theta_1) \\
% \nonumber
% \
% &\quad
% +
% \frac12(\theta-\theta_1)^{\top}
% \widetilde C_{1,S_1}^{\Delta}
% (\theta-\theta_1)
% \nonumber
% \\
% &\quad
% +
% \frac12
% \operatorname{tr}
% \big(
% \widetilde C_{1,S_1}^{\Delta}\Sigma
% \big)
% +
% \mathbb E[R_3(w)].
% \label{eq:expected_risk_gaussian}
% \end{align}

% Step 3: Difference between \(\mathcal P_2\) and \(\mathcal P_1\)

% Apply \eqref{eq:expected_risk_gaussian} to both \((\theta_2,\Sigma_2)\) and \((\theta_1,\Sigma_1)\).
% Since \(\theta_1-\theta_1=0\), we obtain

% \begin{align}
% \Delta^{\mathrm{dist}}_{1\to2}
% &=
% g_{1,S_1}^{\top}\delta_{1\to2}
% +
% \frac12,
% \delta_{1\to2}^{\top}
% \widetilde C_{1,S_1}^{\Delta}
% \delta_{1\to2}
% \nonumber
% \\
% &\quad
% +
% \frac12
% \operatorname{tr}
% \big(
% \widetilde C_{1,S_1}^{\Delta}(\Sigma_2-\Sigma_1)
% \big)
% +
% R_3^{\mathrm{dist}},
% \label{eq:unified_dist_taylor_final}
% \end{align}
% where
% $
% R_3^{\mathrm{dist}}
% =
% \mathbb E_{w\sim\mathcal P_2}[R_3(w)]-
% \mathbb E_{w\sim\mathcal P_1}[R_3(w)].
% $



% Interpretation

% Compared with the pointwise expansion (Part I), the additional term is

% \begin{equation}
% \frac12
% \operatorname{tr}
% \big(
% \widetilde C_{1,S_1}^{\Delta}(\Sigma_2-\Sigma_1)
% \big),
% \label{eq:shape_drift_term_final}
% \end{equation}

% which we call the **shape drift term**.

% It represents:
% \begin{itemize}
%     \item  curvature-weighted covariance shift
%     \item   sensitivity of the old task risk to parameter uncertainty
%     \item  second-order interaction between loss curvature and posterior spread
% \end{itemize}

% In contrast:
% \begin{itemize}
%    \item The first term \(g^\top\delta\) captures first-order interference.
%    \item The second term \(\frac12\delta^\top C\delta\) captures curvature amplification of mean shift.
%    \item The third term \(\frac12\operatorname{tr}(C(\Sigma_2-\Sigma_1))\) captures curvature-weighted uncertainty drift.
% \end{itemize}


% Conceptual hierarchy (clarified)

% This Part II analysis is **not** a KL expansion on the statistical model manifold.
% Instead, it is:

% > A second-order expansion of the empirical risk functional lifted from parameter space to probability measures over parameters.

% In other words:
% \begin{itemize}
% \item Part I: geometry in parameter space.
% \item Part II: geometry in parameter distribution space.
% \end{itemize}
% The trace term appears because
% \[
% \mathbb E[(w-\theta_1)^\top C (w-\theta_1)]
% =
% (\theta-\theta_1)^\top C (\theta-\theta_1)
% +
% \operatorname{tr}(C\Sigma),
% \]
% which introduces covariance sensitivity absent in the pointwise case.




\end{document}